\subsection{局部上同调}

在本节中，我们考虑一个局部 Noether 环 A。回忆（VIII, § 3, 第 3 节，引理 2），A 的定义理想就是 A 的异于 A 且包含 \(\mathfrak{m}_{\mathrm{A}}\) 的一个幂的那些理想，或等价地说，是理想 \(\mathfrak{a} \subset \mathfrak{m}_{\mathrm{A}}\) ，它使得 \(\mathrm{A}/\mathfrak{a}\) 为有限长度。我们用 \(\mathcal{D}\) 表示 \(\mathrm{A}\) 的定义理想组成的集合，并赋以与包含关系相反的序关系；它是右有向的。

设 M 为 A-模。使 A 的每个定义理想 a 对应于分次 A-模 \(\operatorname{Ext}_{\mathrm{A}}(\mathrm{A}/\mathfrak{a},\mathrm{M})\) ；若 a 与 b 是满足 \(a \subset b\) 的定义理想，则用 \(p_{ab}: A/a \to A/b\) 表示典范映射，并考虑 A-线性映射 \(\operatorname{Ext}(p_{ab},1_{\mathrm{M}}): \operatorname{Ext}_{\mathrm{A}}(\mathrm{A}/\mathfrak{b},\mathrm{M}) \longrightarrow \operatorname{Ext}_{\mathrm{A}}(\mathrm{A}/\mathfrak{a},\mathrm{M})\) 。这样我们就得到关于有序集 D 的、由分次 A-模与 0 次分次 A-线性映射组成的归纳系统。我们称 M 的局部上同调模，并记之为 \(\mathrm{H}_{\mathrm{A}}(\mathrm{M})\) ，它就是分次 A-模 \(\lim \operatorname{Ext}_{\mathrm{A}}(\mathrm{A}/\mathfrak{a},\mathrm{M})\) 。

理想 \(\mathfrak{m}_{\mathrm{A}}^{n}\)（\(n \geqslant 1\)）构成 \(\mathcal{D}\) 的一个共尾子集；因此有一个分次模的典范同构 \(\varinjlim_{n} \operatorname{Ext}_{\mathrm{A}}(\mathrm{A}/\mathfrak{m}_{\mathrm{A}}^{n}, M) \longrightarrow \mathrm{H}_{\mathrm{A}}(\mathrm{M})\) 。因而，\(\mathrm{H}_{\mathrm{A}}(\mathrm{M})\) 的每个元素都被 \(\mathfrak{m}_{\mathrm{A}}\) 的一个幂零化。

注 1.—— *设 X 为拓扑空间* \(\operatorname{Spec}(A)\) ，\(\mathcal{O}_{\mathrm{X}}\) 为环的结构层，而 \(\widetilde{\mathbf{M}}\) 为与 M 相伴的 \(\mathcal{O}_{\mathrm{X}}\) -模。分次 A-模 \(\mathrm{H}_{\mathrm{A}}(\mathrm{M})\) 等同于模 \(\mathrm{H}_{\{\mathfrak{m}_{\mathrm{A}}\}}(\mathrm{X},\widetilde{\mathrm{M}})\) ，即以闭点 \(\mathfrak{m}_{\mathrm{A}}\) 为支撑的上同调模（该点在 \(X\) 中是闭的）。

对 A-模的每个同态 \(f: \mathrm{M} \to \mathrm{N}\) ，映射 \(\operatorname{Ext}_{\mathrm{A}}(1_{\mathrm{A}/\mathfrak{a}}, f): \operatorname{Ext}_{\mathrm{A}}(\mathrm{A}/\mathfrak{a}, \mathrm{M}) \longrightarrow \operatorname{Ext}_{\mathrm{A}}(\mathrm{A}/\mathfrak{a}, \mathrm{N})\) 构成一个分次线性映射的归纳系统。取归纳极限，得到一个分次同态 \(\mathrm{H}_{\mathrm{A}}(f): \mathrm{H}_{\mathrm{A}}(\mathrm{M}) \to \mathrm{H}_{\mathrm{A}}(\mathrm{N})\) 。对由 A-模与同态组成的每个序列 \(\mathrm{M} \xrightarrow{f} \mathrm{N} \xrightarrow{g} \mathrm{P}\) ，有 \(\mathrm{H}_{\mathrm{A}}(g \circ f) = \mathrm{H}_{\mathrm{A}}(g) \circ \mathrm{H}_{\mathrm{A}}(f)\) 。设

\[
0 \longrightarrow \mathrm{M} \stackrel {f} {\longrightarrow} \mathrm{N} \stackrel {g} {\longrightarrow} \mathrm{P} \longrightarrow 0\tag{ℰ}
\]

为 A-模的一个正合列。由 A, X, 第 90 页，命题 8，扩张模的连接同态 \(\operatorname{Ext}_{\mathrm{A}}(\mathrm{A}/\mathfrak{a},\mathrm{P}) \longrightarrow \operatorname{Ext}_{\mathrm{A}}(\mathrm{A}/\mathfrak{a},\mathrm{M})\) 构成一个 A-线性映射的归纳系统，这些映射是（升）次数 +1 的分次映射。取归纳极限，导出一个 A-同态 \(\partial(\mathcal{E}): \mathrm{H}_{\mathrm{A}}(\mathrm{P}) \to \mathrm{H}_{\mathrm{A}}(\mathrm{M})\) ，它是次数 +1 的分次映射，并使同态序列

\[
\dots \longrightarrow \mathrm{H} _ {\mathrm{A}} ^ {n - 1} (\mathrm{P}) \xrightarrow {\partial^ {n - 1} (\mathcal {E})} \mathrm{H} _ {\mathrm{A}} ^ {n} (\mathrm{M}) \xrightarrow {\mathrm{H} _ {\mathrm{A}} ^ {n} (f)} \mathrm{H} _ {\mathrm{A}} ^ {n} (\mathrm{N}) \xrightarrow {\mathrm{H} _ {\mathrm{A}} ^ {n} (g)} \mathrm{H} _ {\mathrm{A}} ^ {n} (\mathrm{P}) \xrightarrow {\partial^ {n} (\mathcal {E})} \mathrm{H} _ {\mathrm{A}} ^ {n + 1} (\mathrm{M}) \longrightarrow \dots
\]

设 M 为 A-模。对 A 的每个理想 a，A-模 \(\mathrm{Hom}_{\mathrm{A}}(\mathrm{A}/\mathfrak{a},\mathrm{M})\) 典范地等同于 M 中被 a 零化的元素组成的子模。于是 \(\mathrm{H}_{\mathrm{A}}^{0}(\mathrm{M})\) 等同于 M 中被 \(\mathfrak{m}_{\mathrm{A}}\) 的一个幂零化的元素 m 所组成的子模，亦即满足 \(\mathrm{long}_{\mathrm{A}}(\mathrm{Am}) < +\infty\) 的那些元素组成的子模。特别地，当 M 是 Artin 模时，有 \(\mathrm{H}_{\mathrm{A}}^{0}(\mathrm{M}) = \mathrm{M}\) 。

例.—— 1) 若 M 是内射的，则 A-模 \(\mathrm{H}_{\mathrm{A}}^{i}(\mathrm{M})\) 在 \(i > 0\) 时为零，而在 \(i = 0\) 时是内射的（ \(\S 8, \mathrm{n}^{\circ}2\) ，引理 1, c)）。

\begin{enumerate}
\def\labelenumi{\arabic{enumi})}
\setcounter{enumi}{1}
\item
  若 \(\mathrm{H}_{\mathrm{A}}^{0}(\mathrm{M})=\mathrm{M}\) （例如当 M 是 Artin 模时），则 \(\mathrm{H}_{\mathrm{A}}^{i}(\mathrm{M})\) 在 i > 0 时为零。设 (I, e) 为 M 的一个内射包。I 的子模 \(\mathrm{H}_{\mathrm{A}}^{0}(\mathrm{I})\) 是内射的（例 1）且包含 \(e(\mathrm{M})\) ，故等于 I。令 N 为 e 的余核，并考虑正合列 \(0\rightarrow M\xrightarrow{e}\mathrm{I}\xrightarrow{p}\mathrm{N}\rightarrow0\) 。由于 \(\mathrm{I}=\mathrm{H}_{\mathrm{A}}^{0}(\mathrm{I})\) ，有 \(\mathrm{N}=\mathrm{H}_{\mathrm{A}}^{0}(\mathrm{N})\) ，且同态 \(\mathrm{H}_{\mathrm{A}}^{0}(p)\) 是满射。由于 \(\mathrm{H}_{\mathrm{A}}^{i}(\mathrm{I})\) 在 i > 0 时为零（例 1），\(\mathrm{H}_{\mathrm{A}}^{1}(\mathrm{M})\) 为零，且对 i > 1，\(\mathrm{H}_{\mathrm{A}}^{i}(\mathrm{M})\) 同构于 \(\mathrm{H}_{\mathrm{A}}^{i-1}(\mathrm{N})\) ；对整数 i 作归纳即得结论。
\item
  设 \(\Omega\) 为对偶化 A-模。当 \(i \neq \dim(A)\) 时，有 \(\operatorname{Ext}_{A}^{i}(A/a, \Omega) = 0\) ，其中 \(a\) 是 A 的任意一个定义理想（ \(\S 8\) ，\(n^{\circ} 5\) ，例 3），从而 \(\mathrm{H}_{\mathrm{A}}^{i}(\Omega) = 0\) ；当 \(i = \dim(A)\) 时，A-模 \(\mathrm{H}_{\mathrm{A}}^{i}(\Omega)\) 同构于 \(\varinjlim \operatorname{Ext}_{\mathrm{A}}^{i}(\mathrm{A}/\mathfrak{m}_{\mathrm{A}}^{n}, \Omega)\) ，它是 Matlis A-模（同前引，例 6）。
\item
  设 A 为 Noether 局部整环；用 K 表示它的分式域，并设 A ≠ K。K 是内射 A-模（A, X, 第 18 页，例 1），故模 \(\mathrm{H}_{\mathrm{A}}(K)\) 为零（例 1）。由正合列 0 → A → K → K/A → 0，我们对每个 i 导出同构 \(\mathrm{H}_{\mathrm{A}}^{i}(K/A)\) → \(\mathrm{H}_{\mathrm{A}}^{i+1}(A)\)。
\end{enumerate}

更一般地，对每个无扭 A-模 M 及每个整数 i，我们从正合列

\[
0 \rightarrow \mathrm{M} \rightarrow \mathrm{K} \otimes_ {\mathrm{A}} \mathrm{M} \rightarrow (\mathrm{K/A}) \otimes_ {\mathrm{A}} \mathrm{M} \rightarrow 0
\]

导出同构 \(\mathrm{H}_{\mathrm{A}}^{i}((\mathrm{K}/\mathrm{A})\otimes_{\mathrm{A}}\mathrm{M})\to \mathrm{H}_{\mathrm{A}}^{i + 1}(\mathrm{M})\)

\begin{enumerate}
\def\labelenumi{\arabic{enumi})}
\setcounter{enumi}{4}
\tightlist
\item
  保留上例的假设，并再设 dim(A) = 1。设 N 为扭 A-模；由于 A 的每个异于 A 的非零理想都是定义理想（VIII, § 1, 第 3 节，命题 6, e)），有 \(\mathrm{H}_{\mathrm{A}}^{0}(N) = N\)，从而当 i > 0 时 \(\mathrm{H}_{\mathrm{A}}^{i}(N) = 0\) （例 2）。
\end{enumerate}

设 M 为 A-模；用 T(M) 表示它的扭子模。考虑与正合列相伴随的局部上同调长正合列

\[
0 \to \mathrm{T} (\mathrm{M}) \to \mathrm{M} \to \mathrm{M} / \mathrm{T} (\mathrm{M}) \to 0;
\]

考虑到上述结果，由它导出典范同构 T(M) → \(\mathrm{H}_{\mathrm{A}}^{0}(M)\) 与 \(\mathrm{H}_{\mathrm{A}}^{1}(M)\) → \(\mathrm{H}_{\mathrm{A}}^{1}(M/T(M))\)。由于典范同态 \((K/A) \otimes_{A} M\) → \((K/A) \otimes_{A} (M/T(M))\) 是双射，我们最终得到典范同构

\[
\mathrm{H} _ {\mathrm{A}} ^ {0} (\mathrm{M}) \rightarrow \mathrm{T} (\mathrm{M}), \quad \mathrm{H} _ {\mathrm{A}} ^ {1} (\mathrm{M}) \rightarrow (\mathrm{K} / \mathrm{A}) \otimes_ {\mathrm{A}} \mathrm{M}.
\]

命题 1. 设 A 为 Noether 局部环，M 为有限生成的 A-模。

\begin{enumerate}
\def\labelenumi{\alph{enumi})}
\item
  A-模 \(\mathrm{H}_{\mathrm{A}}(\mathrm{M})\) 是 Artin 模，且在次数 \(>\dim (\mathrm{M})\) 处为零。
\item
  令 \(p = \mathrm{prof}_{\mathrm{A}}(\mathrm{M})\) 。有 \(\mathrm{H}_{\mathrm{A}}^{i}(\mathrm{M}) = 0\) （当 \(i < p\) 时），且 \(\mathrm{H}_{\mathrm{A}}^{p}(\mathrm{M}) \neq 0\) （若 \(M\) 非零）。
\end{enumerate}

我们对 \(\dim(M)\) 作归纳来证明 a)。\(\dim(M) \leqslant 0\) 的情形由上文的例 2 得出。设 \(\dim(M) > 0\) ，先取 \(M\) 形如 \(A/p\) ，其中 \(p\) 是 \(A\) 的异于 \(\mathfrak{m}_A\) 的素理想。设 \(x\) 为 \(\mathfrak{m}_A - p\) 的一个元素；有正合列 \(0 \to M \xrightarrow{x_M} M \longrightarrow M/xM \to 0\) ，其中 \(\dim(M/xM) = \dim(M) - 1\) 。由此导出局部上同调的一个正合列

\[
\mathrm{H} _ {\mathrm{A}} ^ {i - 1} (\mathrm{M} / x \mathrm{M}) \longrightarrow \mathrm{H} _ {\mathrm{A}} ^ {i} (\mathrm{M}) \stackrel {x} {\longrightarrow} \mathrm{H} _ {\mathrm{A}} ^ {i} (\mathrm{M}).
\]

\(\mathrm{H}_{\mathrm{A}}^{i}(\mathrm{M})\) 的每个元素都被 \(\mathfrak{m}_{\mathrm{A}}\) 的一个幂零化；因此，要证明该模是 Artin 模，只需证明 \(\mathrm{H}_{\mathrm{A}}^{i}(\mathrm{M})\) 的底座在 \(\kappa_{A}\) 上是有限维的（§ 8, 第 3 节，引理 3）。由归纳假设，\(\mathrm{H}_{\mathrm{A}}^{i}(\mathrm{M})\) 中以 x 为比的乘法映射的核 N 是 Artin 模；由于 x 属于 \(\mathfrak{m}_{\mathrm{A}}\) ，\(\mathrm{H}_{\mathrm{A}}^{i}(\mathrm{M})\) 的底座等同于 N 的底座，因而是有限维的。若 i > dim(M)，则由归纳假设有 \(\mathrm{H}_{\mathrm{A}}^{i-1}(\mathrm{M}/x\mathrm{M}) = 0\) ，故以 x 为比的乘法映射在 \(\mathrm{H}_{\mathrm{A}}^{i}(\mathrm{M})\) 中是单射；由于 \(\mathrm{H}_{\mathrm{A}}^{i}(\mathrm{M})\) 的每个元素都被 x 的一个幂零化，导出 \(\mathrm{H}_{\mathrm{A}}^{i}(\mathrm{M}) = 0\) ，于是在所考虑的情形 a) 成立。

转到一般情形。A-模 M 容许一个合成列 \((\mathrm{M}_{j})_{0\leqslant j\leqslant n}\) ，使每个商 \(M_{j}/M_{j+1}\) 同构于 \(A/p_{j}\) ，其中 \(p_{j}\) 是 A 的素理想（IV, § 1, 第 4 节，定理 1）。对 n 作归纳，证明 M 满足 a)。n=0 的情形是平凡的。正合列 \(0\to M_{1}\to M\to A/p_{0}\to0\) 给出局部上同调的一个正合列

\[
\mathrm{H} _ {\mathrm{A}} ^ {i} \left(\mathrm{M} _ {1}\right) \longrightarrow \mathrm{H} _ {\mathrm{A}} ^ {i} (\mathrm{M}) \longrightarrow \mathrm{H} _ {\mathrm{A}} ^ {i} \left(\mathrm{A} / \mathfrak {p} _ {0}\right);
\]

由归纳假设，A-模 \(\mathrm{H}_{\mathrm{A}}^{i}(\mathrm{M}_{1})\) 是 Artin 模；由已处理的情形，\(\mathrm{H}_{\mathrm{A}}^{i}(\mathrm{A}/\mathfrak{p}_{0})\) 亦然；于是 \(\mathrm{H}_{\mathrm{A}}^{i}(\mathrm{M})\) 是 Artin 模。若 \(i > \dim(\mathrm{M})\) ，模 \(M_{1}\) 与 \(A/p_{0}\) 的维数 < i；从而由归纳假设与已处理的情形，模 \(\mathrm{H}_{\mathrm{A}}^{i}(\mathrm{M}_{1})\) 与 \(\mathrm{H}_{\mathrm{A}}^{i}(\mathrm{A}/\mathfrak{p}_{0})\) 均为零，这蕴含 \(\mathrm{H}_{\mathrm{A}}^{i}(\mathrm{M}) = 0\) 。

设 M 非零，并对整数 \(p = \text{prof}(M)\) 作归纳来证明 b)。\(p = 0\) 的情形由深度的定义得出。设 \(p > 0\) ，选取一个元素 \(x\) （属于 \(\mathfrak{m}_{\mathrm{A}}\) ），使乘法映射 \(x_{\mathrm{M}}\) 是单射。与上面一样，我们得到局部上同调的一个正合列。

\[
\mathrm{H} _ {\mathrm{A}} ^ {i - 1} (\mathrm{M} / x \mathrm{M}) \longrightarrow \mathrm{H} _ {\mathrm{A}} ^ {i} (\mathrm{M}) \xrightarrow {x} \mathrm{H} _ {\mathrm{A}} ^ {i} (\mathrm{M}).
\]

有 \(\operatorname{prof}(M/xM) = \operatorname{prof}(M) - 1\) （ \(\S 1, n^{\circ} 4\) 命题 7），故由归纳假设得 \(H_{A}^{i-1}(M/xM) = 0\) （对 \(i < p\) ），从而与上面一样推出 \(H_{A}^{i}(M) = 0\) 。特别地，\(H_{A}^{p-1}(M)\) 为零，于是同态 \(H_{A}^{p-1}(M/xM) \longrightarrow H_{A}^{p}(M)\) 是单射；从而由归纳假设，\(H_{A}^{p}(M)\) 非零。

可以证明，当 M 非零时，模 \(H_{A}^{\dim(M)}(M)\) 非零（习题 4；参见第 3 节定理 2 的推论）。

推论. 设 M 为非零的有限生成 Macaulay A-模。A-模 \(\mathrm{H}_{\mathrm{A}}^{i}(\mathrm{M})\) 在 \(i \neq \dim(\mathrm{M})\) 时为零，而在 \(i = \dim(\mathrm{M})\) 时非零。

注 2.—— 对 A 的每个定义理想 a，A-模 \(\operatorname{Ext}_{\mathbf{A}}(\mathbf{A}/\mathfrak{a},\mathbf{M})\) 被 a 零化，且 A/a 等同于 \(\widehat{\mathbf{A}}/\mathfrak{a}\widehat{\mathbf{A}}\) ；因此，分次 A-模 \(\operatorname{Ext}_{\mathbf{A}}(\mathbf{A}/\mathfrak{a},\mathbf{M})\) 等同于 \(\widehat{\mathbf{A}}\otimes_{\mathbf{A}}\operatorname{Ext}_{\mathbf{A}}(\mathbf{A}/\mathfrak{a},\mathbf{M})\) ，从而也等同于 \(\operatorname{Ext}_{\widehat{\mathbf{A}}}(\widehat{\mathbf{A}}/\mathfrak{a}\widehat{\mathbf{A}},\widehat{\mathrm{A}}\otimes_{\mathbf{A}}\mathbf{M})\) （A, X, 第 111 页，命题 10）。诸理想 \(\mathfrak{a}\widehat{\mathbf{A}}\) （其中 \(\mathfrak{a}\in\mathcal{D}\) ）组成的集合包含 \(\mathfrak{m}_{\widehat{\mathbf{A}}}\) 的诸幂，故在 \(\widehat{\mathbf{A}}\) 的定义理想组成的集合中共尾；于是由前述得到分次 A-模的一个同构

\[
\mathrm{H} _ {\mathrm{A}} (\mathrm{M}) \longrightarrow \mathrm{H} _ {\widehat {\mathrm{A}}} (\widehat {\mathrm{A}} \otimes_ {\mathrm{A}} \mathrm{M}).
\]

若 A-模 M 是有限型的，则 A-模 \(\widehat{\mathrm{A}}\otimes_{\mathrm{A}}\mathrm{M}\) 等同于 M 的完备化 \(\widehat{\mathrm{M}}\) （III, § 3, 第 4 节，定理 3），且有一个次数为 0 的分次同构 \(\mathrm{H}_{\mathrm{A}}(\mathrm{M})\to \mathrm{H}_{\widehat{\mathrm{A}}}(\widehat{\mathrm{M}})\) 。

\subsection{Macaulay 环上的局部上同调}

在本节中，我们假设 A 是局部 Macaulay 环；令 \(\dim(A) = d\) 。

由 \(\mathfrak{m}_{\mathrm{A}}\) 中对 A 完全割、且长度为 \(d = \dim (\mathrm{A})\) 的元素序列所生成的诸理想，构成子集 \(\mathcal{D}_{cs}\) ，它在 A 的定义理想组成的集合 \(\mathcal{D}\) 中共尾。事实上，设 \((x_{1},\ldots ,x_{d})\) 为 \(\mathfrak{m}_{\mathrm{A}}\) 中对 A 完全割的序列（§ 2, 第 3 节，命题 3）；对每个整数 \(n\) ，序列 \((x_1^n,\dots ,x_d^n)\) 对 A 是完全割的（A, X, 第 158 页，命题 6, c)），且生成一个定义理想（VIII, § 3, 第 2 节，命题 3 的推论及定理 1），该理想含于 \(\mathfrak{m}_{\mathrm{A}}^n\) 。

设 \(\mathfrak{a} \in \mathcal{D}_{cs}\) ，设 \(\pi: L \to A/\mathfrak{a}\) 为一个有限型的自由分解，它在次数 \(>d\) 处为零（例如与生成 \(\mathfrak{a}\) 的、对 A 完全割的序列相伴随的 Koszul 复形）。考虑 \(L^{*} = \operatorname{Homgr}_{A}(L,A)\) ，即 \(L\) 的对偶；由于 A 的深度等于 \(d\) ，有 \(\operatorname{Ext}_{A}^{i}(A/\mathfrak{a}, A) = 0\) （对 \(i < d\) ）（ \(\S\) 1, \(n^{\circ}\) 1, 命题 2 的推论 2）。由于 \(L^{*}\) 的长度 \(\leqslant d\) ，我们导出：\(H^{i}(L^{*})\) 在 \(i \neq d\) 时为零，且 \(H^{d}(L^{*})\) 等同于 \(\operatorname{Ext}_{A}^{d}(A/\mathfrak{a}, A)\) （A, X, 第 100 页，定理 1）。于是，我们有一个同态

\[
\pi^ {*}: \mathrm{L} ^ {*} (- d) \longrightarrow \operatorname{Ext} _ {\mathrm{A}} ^ {d} (\mathrm{A} / \mathfrak {a}, \mathrm{A})
\]

它定义了 \(\operatorname{Ext}_{\mathrm{A}}^{d}(\mathrm{A}/\mathfrak{a},\mathrm{A})\) 的一个有限型自由分解。

设 M 为 A-模；考虑典范同构（同前引）

\[
\varphi (\mathrm{L}, \mathrm{M}): \mathrm{H} (\operatorname{Homgr} _ {\mathrm{A}} (\mathrm{L}, \mathrm{M})) \rightarrow \operatorname{Ext} _ {\mathrm{A}} (\mathrm{A} / \mathfrak {a}, \mathrm{M})
\]

\[
\psi (\mathrm{M}, \mathrm{L} ^ {*} (- d)): \mathrm{Tor} ^ {\mathrm{A}} (\mathrm{M}, \mathrm{Ext} _ {\mathrm{A}} ^ {d} (\mathrm{A} / \mathfrak {a}, \mathrm{A})) \to \mathrm{H} (\mathrm{M} \otimes_ {\mathrm{A}} \mathrm{L} ^ {*}) (- d).
\]

由于复形 L 是有限型自由的，复形的典范态射 \(M \otimes_{A} L^{*} \to \text{Homgr}_{A}(L, M)\) 是一个同构；由此导出分次 A-模的一个同构 \(H(M \otimes_{A} L^{*}) \to H(\text{Homgr}_{A}(L, M))\) 。合成上述诸同构，我们得到分次 A-模的一个同构，它称为典范同构。

\[
\tau (\mathrm{L}, \mathrm{M}): \operatorname{Tor} ^ {\mathrm{A}} (\mathrm{M}, \operatorname{Ext} _ {\mathrm{A}} ^ {d} (\mathrm{A} / \mathfrak {a}, \mathrm{A})) (d) \longrightarrow \operatorname{Ext} _ {\mathrm{A}} (\mathrm{A} / \mathfrak {a}, \mathrm{M})
\]

它对每个整数 i 诱导一个同构

\[
\tau^ {i} (\mathrm{L}, \mathrm{M}): \operatorname{Tor} _ {d - i} ^ {\mathrm{A}} (\mathrm{M}, \operatorname{Ext} _ {\mathrm{A}} ^ {d} (\mathrm{A} / \mathfrak {a}, \mathrm{A})) \longrightarrow \operatorname{Ext} _ {\mathrm{A}} ^ {i} (\mathrm{A} / \mathfrak {a}, \mathrm{M}).
\]

对 \(M = A\) ，\(\tau^d(L, A)\) 是从 \(A \otimes_A \operatorname{Ext}_A^d(\mathrm{A} / \mathfrak{a}, A)\) 到 \(\operatorname{Ext}_\mathrm{A}^d(A / \mathfrak{a}, A)\) 上的典范同构。

设 \(\mathfrak{b}\) 为 \(\mathcal{D}_{cs}\) 中含于 \(\mathfrak{a}\) 的一个理想。设 \(\rho: R \to A/\mathfrak{b}\) 为长度 \(\leqslant d\) 的一个有限自由分解，并设 \(p_{\mathfrak{ab}}: A/\mathfrak{b} \to A/\mathfrak{a}\) 为典范满射。由 A, X, 第 49 页，命题 3，存在复形的一个态射 \(\mathrm{P}_{\mathrm{LR}}: R \to L\) ，使得 \(\pi \circ \mathrm{P}_{\mathrm{LR}} = p_{\mathfrak{ab}} \circ \rho\) 。由 A, X, 第 103 页，命题 2，我们有一个交换图

\[
\begin{array}{c c c} \operatorname{Tor} ^ {A} (M, \operatorname{Ext} _ {A} ^ {d} (A / \mathfrak {a}, A)) (d) & \xrightarrow {\operatorname{Tor} (1 _ {M} , \operatorname{Ext} ^ {d} (p _ {a b} , 1 _ {\mathrm{A}}))} & \operatorname{Tor} ^ {\mathrm{A}} (M, \operatorname{Ext} _ {A} ^ {d} (A / \mathfrak {b}, A)) (d) \\ \psi (L ^ {*} (- d), M) \Bigg \downarrow & & \Bigg \downarrow \psi (R ^ {*} (- d), M) \\ H (M \otimes_ {A} L ^ {*}) & \xrightarrow {H (1 _ {M} \otimes^ {t} P _ {L R})} & H (M \otimes_ {A} R ^ {*}) \\ \Bigg \downarrow & & \Bigg \downarrow \\ H (\operatorname{Homgr} _ {\mathrm{A}} (L, M)) & \xrightarrow {H (\operatorname{Homgr} (P _ {L R} , M))} & H (\operatorname{Homgr} _ {A} (R, M)) \\ \varphi (L, M) \Bigg \downarrow & & \Bigg \downarrow \varphi (R, M) \\ \operatorname{Ext} _ {\mathrm{A}} (A / \mathfrak {a}, M) & \xrightarrow {\operatorname{Ext} (p _ {a b} , 1 _ {M})} & \operatorname{Ext} _ {A} (A / \mathfrak {b}, M) \end{array}
\]

首先，取 \(\mathfrak{a} = \mathfrak{b}\) ，即可看出同构 \(\tau(\mathrm{L},\mathrm{M})\) 不依赖于 A/ \(\mathfrak{a}\) 的分解 L 的选取；我们把它记作 \(\tau_{\mathfrak{a}}(\mathrm{M})\) 。由此进而推出，诸 \(\tau_{\mathfrak{a}}(\mathrm{M})\) （其中 \(\mathfrak{a} \in \mathcal{D}_{cs}\) ）构成同构的一个归纳系统。取归纳极限，并注意到 A, X, 第 70 页，命题 8，对每个整数 \(i\) 得到 A-模的一个同构

\[
\tau^ {i} (\mathrm{M}): \operatorname{Tor} _ {d - i} ^ {\mathrm{A}} (\mathrm{M}, \mathrm{H} _ {\mathrm{A}} ^ {d} (\mathrm{A})) \longrightarrow \mathrm{H} _ {\mathrm{A}} ^ {i} (\mathrm{M}).
\]

对 M = A，\(\tau^{d}(A)\) 是从 \(A \otimes_{A} H_{A}^{d}(A)\) 到 \(H_{A}^{d}(A)\) 上的典范同构。

注.—— 1) 设 \(f: \mathrm{M} \to \mathrm{N}\) 为 A-模的一个同态。利用 A, X, 第 103 页，命题 2，可以证明下列诸图是交换的：

\[
\begin{array}{c c c} \operatorname{Tor} _ {d - i} ^ {\mathrm{A}} (\mathrm{M}, \mathrm{H} _ {\mathrm{A}} ^ {d} (\mathrm{A})) & \xrightarrow {\tau^ {i} (\mathrm{M})} & \mathrm{H} _ {\mathrm{A}} ^ {i} (\mathrm{M}) \\ \operatorname{Tor} _ {d - i} (f, 1) \Bigg \downarrow & & \Bigg \downarrow \mathrm{H} ^ {i} (f) \\ \operatorname{Tor} _ {d - i} ^ {\mathrm{A}} (\mathrm{N}, \mathrm{H} _ {\mathrm{A}} ^ {d} (\mathrm{A})) & \xrightarrow {\tau^ {i} (\mathrm{N})} & \mathrm{H} _ {\mathrm{A}} ^ {i} (\mathrm{N}) \quad . \end{array}
\]

\begin{enumerate}
\def\labelenumi{\arabic{enumi})}
\setcounter{enumi}{1}
\tightlist
\item
  设
\end{enumerate}

\[
0 \to \mathrm{M} \to \mathrm{N} \to \mathrm{P} \to 0\tag{ℰ}
\]

为 A-模的一个正合列。利用 A, X, 第 104 页，命题 3 以及第 106 页，命题 4，可以证明下列诸图是交换的：

\[
\begin{array}{c c c} \operatorname{Tor} _ {d - i} ^ {\mathrm{A}} (\mathrm{P}, \mathrm{H} _ {\mathrm{A}} ^ {d} (\mathrm{A})) & \xrightarrow {\tau^ {i} (\mathrm{P})} & \mathrm{H} _ {\mathrm{A}} ^ {i} (\mathrm{P}) \\ \partial_ {d - i} (\mathscr {E}, \mathrm{H} _ {\mathrm{A}} ^ {d} (\mathrm{A})) \Bigg \downarrow & & \Bigg \downarrow \partial^ {i} (\mathscr {E}) \\ \operatorname{Tor} _ {d - i - 1} ^ {\mathrm{A}} (\mathrm{M}, \mathrm{H} _ {\mathrm{A}} ^ {d} (\mathrm{A})) & \xrightarrow {\tau^ {i + 1} (\mathrm{M})} & \mathrm{H} _ {\mathrm{A}} ^ {i + 1} (\mathrm{M}) \end{array} .
\]

\begin{enumerate}
\def\labelenumi{\arabic{enumi})}
\setcounter{enumi}{2}
\tightlist
\item
  考虑同构
\end{enumerate}

\[
\tau^ {d} (\mathrm{M}): \mathrm{M} \otimes_ {\mathrm{A}} \mathrm{H} _ {\mathrm{A}} ^ {d} (\mathrm{A}) \longrightarrow \mathrm{H} _ {\mathrm{A}} ^ {d} (\mathrm{M}).
\]

对 \(x \in \mathrm{M}\) ，用 \(f_x\) 表示从 A 到 M 的映射 \(a \mapsto ax\) 。对每个 \(u \in \mathrm{H}_{\mathrm{A}}^{d}(\mathrm{A})\) ，有

\[
\tau^ {d} (\mathrm{M}) (x \otimes u) = \mathrm{H} _ {\mathrm{A}} ^ {d} (f _ {x}) (u).
\]

这确实由把注 1 应用于同态 \(f_{x}:A\to M\) 而得。

\subsection{Macaulay 环上的 Grothendieck 对偶}

我们继续假设 A 是维数为 \(d\) 的局部 Macaulay 环。设 \(\Omega\) 为一个对偶化 A-模。A-模 \(\mathrm{H}_{\mathrm{A}}^{d}(\Omega)\) 是一个 Matlis 模（ \(\S 8\) , 第 5 节，例 6）；按照 \(\S 8\) 的记号，对每个 A-模 M，令 \(\mathrm{D(M)} = \mathrm{Hom}_{\mathrm{A}}(\mathrm{M}, \mathrm{H}_{\mathrm{A}}^{d}(\Omega))\) 。

考虑同构 \(\tau^{d}(\Omega):\Omega\otimes\mathrm{H}_{\mathrm{A}}^{d}(\mathrm{A})\to\mathrm{H}_{\mathrm{A}}^{d}(\Omega)\) （第 2 节）。由此导出一个同态 \(\omega:\mathrm{H}_{\mathrm{A}}^{d}(\mathrm{A})\to\mathrm{Hom}_{\mathrm{A}}(\Omega,\mathrm{H}_{\mathrm{A}}^{d}(\Omega))\) ，它把每个元素 \(u\) （属于 \(\mathrm{H}_{\mathrm{A}}^{d}(\mathrm{A})\) ）对应到同态 \(x\mapsto\mathrm{H}_{\mathrm{A}}^{d}(f_{x})(u)\) （见上面的注 3）。

命题 2.—— 设 A 为维数 d 的局部 Macaulay 环，并设 \(\Omega\) 为一个对偶化 A-模。同态 \(\omega: H_{A}^{d}(A) \to D(\Omega)\) 是双射。

同态 \(\omega\) 是映射的归纳系统 \((\omega_{\mathfrak{a}})_{\mathfrak{a}\in\mathcal{D}_{cs}}\) 的极限，其中

\[
\omega_ {\mathfrak {a}}: \operatorname{Ext} _ {\mathrm{A}} ^ {d} (\mathrm{A} / \mathfrak {a}, \mathrm{A}) \longrightarrow \operatorname{Hom} _ {\mathrm{A}} (\Omega , \operatorname{Ext} _ {\mathrm{A}} ^ {d} (\mathrm{A} / \mathfrak {a}, \Omega))
\]

它把每个元素 \(u\) （属于 \(\operatorname{Ext}_{\mathrm{A}}^{d}(\mathrm{A}/\mathfrak{a},\mathrm{A})\) ）对应到映射 \(x\mapsto f_{x}\circ u\) （A, X, 第 114 页）。因此只需证明每个映射 \(\omega_{\mathfrak{a}}\) 都是双射。

设 \(\mathfrak{a}\) 为 \(\mathcal{D}_{cs}\) 中由对 A 完全割的序列 \(\mathbf{x} = (x_1, \ldots, x_d)\) 所生成的一个理想。Koszul 复形 \(\mathbf{K}^{\bullet}(\mathbf{x}, \mathrm{A})\) 给出 A/ \(\mathfrak{a}\) 的一个投射分解；对每个 A-模 M，A-模 \(\mathrm{H}^d\big(\mathrm{Homgr}_\mathrm{A}(\mathbf{K}^{\bullet}(\mathbf{x}, \mathrm{A}), \mathrm{M})\big)\) 典范地等同于 M/ \(\mathfrak{a}\mathrm{M}\) （A, X, 第 155 页）。由此导出一个同构（A, X, 第 100 页）

\[
\varphi_ {\mathrm{M}}: \mathrm{M} / \mathfrak {a} \mathrm{M} \longrightarrow \operatorname{Ext} ^ {d} (\mathrm{A} / \mathfrak {a}, \mathrm{M}).
\]

设 \(x \in \Omega\) 。注意到同前引第 103 页，命题 2，我们有一个交换图

\[
\begin{array}{c c c} \mathrm{A/a} & \xrightarrow {\varphi_ {\mathrm{A}}} & \operatorname{Ext} _ {\mathrm{A}} ^ {d} (\mathrm{A/a}, \mathrm{A}) \\ \bar {f} _ {x} \Bigg \downarrow & & \Bigg \downarrow \operatorname{Ext} (1 _ {\mathrm{A} / a}, f _ {x}) \\ \Omega / a \Omega & \xrightarrow {\varphi_ {\Omega}} & \operatorname{Ext} _ {\mathrm{A}} ^ {d} (\mathrm{A/a}, \Omega) \end{array}
\]

其中 \(\bar{f}_x\) 是由 \(f_x\) 通过取商导出的同态。由此可知，若对每个 A-模 M，把 \(\mathrm{Ext}_{\mathrm{A}}^{d}(\mathrm{A}/\mathfrak{a},\mathrm{M})\) 与 M/ \(\mathfrak{a}\) M （借助 \(\varphi_{\mathrm{M}}\) ）等同起来，则同态 \(\omega_{\mathfrak{a}}\) 就等同于从 A/ \(\mathfrak{a}\) 到 \(\mathrm{Hom}_{\mathrm{A}}(\Omega,\Omega/\mathfrak{a}\Omega)\) 的、把 1 映为典范满射的那个 A-线性映射，也就是典范映射 A/ \(\mathfrak{a}\longrightarrow\mathrm{End}_{\mathrm{A}/\mathfrak{a}}(\Omega/\mathfrak{a}\Omega)\) 。但由于 A/ \(\mathfrak{a}\) -模 \(\Omega/\mathfrak{a}\Omega\) 是对偶化的（\(\S\) 9, 第 2 节，命题 4），这个映射是双射（ \(\S\) 9, 第 4 节，命题 6），这就证明了本命题。

借助同构 \(\widehat{\alpha}_{\Omega}\) （§ 8, 第 3 节，定理 2, b)），把 Matlis 双对偶 D(D(Ω)) 与 Ω 等同起来。

推论.—— 同态 D(ω): \(\widehat{\Omega} \rightarrow \mathrm{D}(\mathrm{H}_{\mathrm{A}}^{d}(\mathrm{A}))\) 是一个同构。

设 M 为一个 A-模，并设 i 为一个整数。考虑典范同态（§ 8, 第 7 节）

\[
\rho_ {d - i} (\mathrm{M}, \Omega): \mathrm{Tor} _ {d - i} ^ {\mathrm{A}} (\mathrm{M}, \mathrm{D} (\Omega)) \longrightarrow \mathrm{D} \big (\mathrm{Ext} _ {\mathrm{A}} ^ {d - i} (\mathrm{M}, \Omega) \big)
\]

\[
\theta^ {d - i} (\mathrm{M}, \mathrm{H} _ {\mathrm{A}} ^ {d} (\mathrm{A})): \operatorname{Ext} _ {\mathrm{A}} ^ {d - i} \left(\mathrm{M}, \mathrm{D} \left(\mathrm{H} _ {\mathrm{A}} ^ {d} (\mathrm{A})\right)\right) \longrightarrow \mathrm{D} \left(\operatorname{Tor} _ {d - i} ^ {\mathrm{A}} \left(\mathrm{M}, \mathrm{H} _ {\mathrm{A}} ^ {d} (\mathrm{A})\right)\right).
\]

利用同构 \(\omega : \mathrm{H}_{\mathrm{A}}^{d}(\mathrm{A}) \to \mathrm{D}(\Omega)\) 、\(\mathrm{D}(\omega) : \widehat{\Omega} \to \mathrm{D}(\mathrm{H}_{\mathrm{A}}^{d}(\mathrm{A}))\) （命题 2 的推论 1）以及 \(\tau^{i}(\mathrm{M}) : \mathrm{Tor}_{d-i}^{\mathrm{A}}(\mathrm{M}, \mathrm{H}_{\mathrm{A}}^{d}(\mathrm{A})) \longrightarrow \mathrm{H}_{\mathrm{A}}^{i}(\mathrm{M})\) （ \(n^{\circ}\) 2），我们导出 A-模的典范同态

\[
\gamma^ {i} (\mathrm{M}): \mathrm{H} _ {\mathrm{A}} ^ {i} (\mathrm{M}) \longrightarrow \mathrm{D} \left(\mathrm{Ext} _ {\mathrm{A}} ^ {d - i} (\mathrm{M}, \Omega)\right)
\]

\[
\delta^ {i} (\mathrm{M}): \operatorname{Ext} _ {\mathrm{A}} ^ {d - i} (\mathrm{M}, \widehat {\Omega}) \longrightarrow \mathrm{D} \left(\mathrm{H} _ {\mathrm{A}} ^ {i} (\mathrm{M})\right).
\]

定理 1（Grothendieck 对偶）. 设 A 为维数 d 的局部 Macaulay 环，并设 Ω 为一个对偶化 A-模。

\begin{enumerate}
\def\labelenumi{\alph{enumi})}
\item
  A-模 \(\mathrm{H}_{\mathrm{A}}^{d}(\Omega)\) 是一个 Matlis 模；对每个 A-模 P，用 D(P) 表示 Matlis 对偶 \(\mathrm{Hom}_{\mathrm{A}}(\mathrm{P},\mathrm{H}_{\mathrm{A}}^{d}(\Omega))\) 。
\item
  对每个有限型 \(\mathrm{A}\) -模 M 以及任意整数 \(i\) ，典范同态
\end{enumerate}

\[
\gamma^ {i} (\mathrm{M}): \mathrm{H} _ {\mathrm{A}} ^ {i} (\mathrm{M}) \longrightarrow \mathrm{D} \left(\operatorname{Ext} _ {\mathrm{A}} ^ {d - i} (\mathrm{M}, \Omega)\right)
\]

是 Artin A-模之间的一个同构。

\begin{enumerate}
\def\labelenumi{\alph{enumi})}
\setcounter{enumi}{2}
\tightlist
\item
  对每个 A-模 M 和每个整数 i，典范同态
\end{enumerate}

\[
\delta^ {i} (\mathrm{M}): \operatorname{Ext} _ {\mathrm{A}} ^ {d - i} (\mathrm{M}, \widehat {\Omega}) \longrightarrow \mathrm{D} \left(\mathrm{H} _ {\mathrm{A}} ^ {i} (\mathrm{M})\right)
\]

是 \(\widehat{\mathrm{A}}\) -模之间的一个同构。

这由 § 8, 第 7 节，命题 6 得出。

推论.—— 设 A 为一个局部 Macaulay 环，M 为一个维数为 e 的非零有限生成 A-模。A-模 \(\mathrm{H}_{\mathrm{A}}^{e}(\mathrm{M})\) 非零。

根据第 1 节的注 2，我们可以假设局部环 A 是完备的。在此情形 A 具有一个对偶化模 Ω（§ 9, 第 3 节，命题 6 的推论 3）；若 \(\mathrm{H}_{\mathrm{A}}^{e}(M)\) 为零，则它的 Matlis 对偶 \(\operatorname{Ext}_{\mathrm{A}}^{d-e}(M, \Omega)\) 也为零，这与 § 9, 第 1 节，命题 3, b) 相矛盾。

注.—— 1) 当 A-模 M 是有限生成的时，\(\widehat{\mathrm{A}}\) -模 \(\operatorname{Ext}_{\mathrm{A}}^{d-i}(\mathrm{M}, \widehat{\Omega})\) 等同于 \(\widehat{\mathrm{A}} \otimes_{\mathrm{A}} \operatorname{Ext}_{\mathrm{A}}^{d-i}(\mathrm{M}, \Omega)\) （A, X, 第 108 页，命题 7, c)），且 \(\delta^{i}(\mathrm{M})\) 也可由 \(\mathrm{D}(\gamma^{i}(\mathrm{M}))\) 与双对偶性同构复合而得到。

\begin{enumerate}
\def\labelenumi{\arabic{enumi})}
\setcounter{enumi}{1}
\tightlist
\item
  设 \(u: \mathrm{M} \to \mathrm{M}'\) 为 A-模的一个同态。由第 2 节的注 1 以及 § 8, 第 7 节的注，下列诸图是交换的：
\end{enumerate}

\[
\begin{array}{c c c} \mathrm{H} _ {\mathrm{A}} ^ {i} (\mathrm{M}) & \xrightarrow {\gamma^ {i} (\mathrm{M})} & \mathrm{D} (\mathrm{Ext} _ {\mathrm{A}} ^ {d - i} (\mathrm{M}, \Omega)) \\ \mathbb {1} _ {\mathrm{A}} ^ {i} (u) \Bigg \downarrow & & \Bigg \downarrow \mathrm{D} (\mathrm{Ext} (u, 1)) \\ \mathrm{H} _ {\mathrm{A}} ^ {i} (\mathrm{M} ^ {\prime}) & \xrightarrow {\gamma^ {i} (\mathrm{M} ^ {\prime})} & \mathrm{D} (\mathrm{Ext} _ {\mathrm{A}} ^ {d - i} (\mathrm{M} ^ {\prime}, \Omega)) \end{array}
\]

\[
\begin{array}{c c c} \operatorname{Ext} _ {\mathrm{A}} ^ {d - i} (\mathrm{M} ^ {\prime}, \widehat {\Omega}) & \xrightarrow {\delta^ {i} (\mathrm{M} ^ {\prime})} & \mathrm{D} (\mathrm{H} _ {\mathrm{A}} ^ {i} (\mathrm{M} ^ {\prime})) \\ \operatorname{Ext} (u, 1) \Bigg \downarrow & & \Bigg \downarrow \mathrm{D} (\mathrm{H} _ {\mathrm{A}} ^ {i} (u)) \\ \operatorname{Ext} _ {\mathrm{A}} ^ {d - i} (\mathrm{M}, \widehat {\Omega}) & \xrightarrow {\delta^ {i} (\mathrm{M})} & \mathrm{D} (\mathrm{H} _ {\mathrm{A}} ^ {i} (\mathrm{M})) \end{array} .
\]

\begin{enumerate}
\def\labelenumi{\arabic{enumi})}
\setcounter{enumi}{2}
\tightlist
\item
  设
\end{enumerate}

\[
0 \rightarrow \mathrm{M} ^ {\prime} \rightarrow \mathrm{M} \rightarrow \mathrm{M} ^ {\prime \prime} \rightarrow 0\tag{ℰ}
\]

为 A-模的一个正合列。由第 2 节的注 2 以及 § 8, 第 7 节的注，下列诸图是交换的：

\[
\begin{array}{c c c} \mathrm{H} _ {\mathrm{A}} ^ {i - 1} (\mathrm{M} ^ {\prime \prime}) & \xrightarrow {\gamma^ {i - 1} (\mathrm{M} ^ {\prime \prime})} & \mathrm{D} (\mathrm{Ext} _ {\mathrm{A}} ^ {d - i + 1} (\mathrm{M} ^ {\prime \prime}, \Omega)) \\ \partial^ {i - 1} (\mathscr {E}) \Bigg \downarrow & & \Bigg \downarrow (- 1) ^ {d - i + 1} \mathrm{D} (\delta^ {d - i} (\mathscr {E}, \Omega)) \\ \mathrm{H} _ {\mathrm{A}} ^ {i} (\mathrm{M} ^ {\prime}) & \xrightarrow {\gamma^ {i} (\mathrm{M} ^ {\prime})} & \mathrm{D} (\mathrm{Ext} _ {\mathrm{A}} ^ {d - i} (\mathrm{M} ^ {\prime}, \Omega)) \end{array}
\]

\[
\begin{array}{c c c} \operatorname{Ext} _ {\mathrm{A}} ^ {d - i} (\mathrm{M} ^ {\prime}, \widehat {\Omega}) & \xrightarrow {\delta^ {i} (\mathrm{M} ^ {\prime})} & \mathrm{D} (\mathrm{H} _ {\mathrm{A}} ^ {i} (\mathrm{M} ^ {\prime})) \\ \delta^ {d - i} (\mathscr {E}, \widehat {\Omega}) \Bigg \downarrow & & \Bigg \downarrow (- 1) ^ {d - i + 1} \mathrm{D} (\partial^ {i - 1} (\mathscr {E})) \\ \operatorname{Ext} _ {\mathrm{A}} ^ {d - i + 1} (\mathrm{M} ^ {\prime \prime}, \widehat {\Omega}) & \xrightarrow {\delta^ {i - 1} (\mathrm{M} ^ {\prime \prime})} & \mathrm{D} (\mathrm{H} _ {\mathrm{A}} ^ {i - 1} (\mathrm{M} ^ {\prime \prime})) \quad . \end{array}
\]

例.—— 设 A 为维数 1 的 Noether 局部整环；用 K 表示它的分式域。设 Ω 为一个对偶化 A-模，并设 M 为一个有限生成 A-模。A-模 \(\mathrm{H}_{\mathrm{A}}^{0}(M)\) 与 \(\mathrm{H}_{\mathrm{A}}^{1}(M)\) 分别典范地等同于 T(M) 与 \((K/A) \otimes_{A} M\) （第 1 节，例 5）。在这些等同之下，对偶同构

\[
\gamma^ {0} (\mathrm{M}): \mathrm{T} (\mathrm{M}) \longrightarrow \mathrm{D} \left(\operatorname{Ext} _ {\mathrm{A}} ^ {1} (\mathrm{M}, \Omega)\right), \quad \gamma^ {1} (\mathrm{M}): (\mathrm{K} / \mathrm{A}) \otimes_ {\mathrm{A}} \mathrm{M} \longrightarrow \mathrm{D} \left(\operatorname{Hom} _ {\mathrm{A}} (\mathrm{M}, \Omega)\right)
\]

（定理 1）就是 § 9, 第 6 节中定义的那些同构。

\specialchapter{习题}

\exercisesection{深度}

\begin{enumerate}
\def\labelenumi{\arabic{enumi})}
\item
  设 A 为环，J 为 A 的理想，M 为 A-模。证明：为使 \(\operatorname{prof}_{\mathrm{A}}(J;M)\geqslant2\) 成立，必要且充分的是，由 J 到 A 中的典范单射所导出的映射 \(M\to\operatorname{Hom}_{\mathrm{A}}(J,M)\) 是一个同构。
\item
  设 A 为一个高为 1 的赋值环，且不是 Noether 的，m 为其极大理想，n 为 A 的一个主理想，异于 (0) 且异于 A。证明我们有 \(\operatorname{prof}_{\mathrm{A}}(\mathfrak{n};\mathrm{A})=1\) 以及 \(\operatorname{prof}_{\mathrm{A}}(\mathfrak{m};\mathrm{A})\geqslant2\) ，尽管 \(\mathrm{V}(\mathfrak{m})=\mathrm{V}(\mathfrak{n})\) （注意 \(\operatorname{Ext}_{\mathrm{A}}^{1}(\mathrm{A}/\mathfrak{m},\mathrm{A}/\mathfrak{n})\) 同构于 \(\operatorname{Hom}_{\mathrm{A}}(\mathrm{A}/\mathfrak{m},\mathrm{A})\) 。）
\item
  设 A 为局部环 \(k[[X,Y]]\) ，并设 M 为诸 A-模 \(\mathrm{A}/\mathfrak{a}\) 的直和，其中 \(\mathfrak{a}\) 取遍 A 的非零主理想组成的集合。证明有 \(\mathrm{prof}_{\mathrm{A}}(\mathrm{M}) = 1\) ，但在 \(\mathfrak{m}_{\mathrm{A}}\) 中不存在 M-正则的元素。
\item
  设 A 为环，J 为 A 的一个有限生成理想，M 为 A-模，P 为平坦 A-模。证明有 \(\operatorname{prof}_{\mathrm{A}}(\mathrm{J};\mathrm{P}\otimes_{\mathrm{A}}\mathrm{M})\geqslant\operatorname{prof}_{\mathrm{A}}(\mathrm{J};\mathrm{M})\) ，且当 P 忠实平坦时等号成立。
\item
  设 \(\mathrm{A}\) 为 Noether 环，M 为一个非零的有限生成 A-模，J 为 \(\mathrm{A}\) 的理想，\(\mathbf{x} = (x_1, \ldots, x_n)\) 为 J 的一个生成组。证明有 \(\mathrm{H}^i(\mathbf{x}, \mathrm{M}) \neq 0\) （对 \(\operatorname{prof}_{\mathrm{A}}(\mathrm{J}; \mathrm{M}) \leqslant i \leqslant n\) ）（化为局部情形，并利用 A, X, 第 157 页，推论 2）。
\item
  设 A 为 Noether 局部环，并设 P 为由平坦 A-模组成的一个长度有限为 \(\ell\) 的复形，使得 \(\operatorname{Supp}(\mathrm{H}(\mathrm{P})) = \{\mathfrak{m}_{\mathrm{A}}\}\) 。证明每个使得 \(\kappa_{\mathrm{A}} \otimes_{\mathrm{A}} \mathrm{M} \neq 0\) 的 A-模 M 的深度都 \(\leqslant \ell\) （把命题 3 应用于复形 \(\mathrm{P} \otimes_{\mathrm{A}} \mathrm{M}\) ，注意到这个复形不是正合的，并利用习题 4）。
\item
  设 A 为局部环，M 为有限生成 A-模，F 为 \(\operatorname{Spec}(A)\) 的一个闭子集。证明不等式 \(\operatorname{prof}_{\mathbf{F}}(\mathbf{M}) \geqslant \operatorname{prof}(\mathbf{M}) - \dim(\mathbf{F})\) 。
\item
  设 A 为 Noether 环，并设 M 为一个有限生成 A-模。我们称 M 满足性质 \((S_{k})\) ，如果对 A 的每个素理想 p 有 \(\operatorname{prof}_{A_{\mathfrak{p}}}(M_{\mathfrak{p}}) \geqslant \inf(k, \dim_{A_{\mathfrak{p}}}(M_{\mathfrak{p}}))\) 。
\end{enumerate}

\begin{enumerate}
\def\labelenumi{\alph{enumi})}
\item
  每个模都满足 \((\mathrm{S}_{k})\) （对 \(k \leqslant 0\) ）；说一个模满足 \((\mathrm{S}_{1})\) ，意味着它没有嵌入的相伴素理想。
\item
  设 \(0 \to M' \to M \to M'' \to 0\) 为有限生成 A-模的一个短正合列，并设 \(k\) 为整数。我们假设 \(M\) 满足 \((S_k)\) ，\(M''\) 满足 \((S_{k-1})\) ，且当 \(k \geqslant 2\) 时有 \(\operatorname{Ass}(M'') \subset \operatorname{Ass}(M)\) 。证明 \(M'\) 满足 \((S_k)\) （利用第 7 节的推论 2）。
\item
  c) 设 \((L,p)\) 为 M 的一个由有限生成自由 A-模组成的分解。若 A 满足性质 \((S_{k})\) ，则 A-模 \(B_{i}(L)\) （对 \(i \geqslant 0\) ）满足 \((S_{h})\) ，其中 \(h = \inf(k,i+1)\) 。 d) 若 A 满足 \((S_{2})\) ，则每个有限生成 A-模的对偶都满足 \((S_{2})\) （应用 c）。 e) 设 J 为 A 的一个理想，由一个 M-正则序列 \((x_{1},\ldots,x_{r})\) 生成。若 M 满足 \((S_{k})\) ，则 A-模 M/JM 满足 \((S_{k-r})\) 。
\end{enumerate}

\(f)\) 设 B 为 Noether A-代数，N 为一个有限生成 B-模且是忠实平坦的 A-模，并设 \(k\) 为整数。若 B-模 \(\mathrm{M} \otimes_{\mathrm{A}} \mathrm{N}\) 满足性质 \((\mathrm{S}_k)\) ，则 A-模 M 亦满足。若 M 满足 \((\mathrm{S}_k)\) ，且 \((\kappa(\mathfrak{p}) \otimes_{\mathrm{A}} \mathrm{B})\) -模 \(\kappa(\mathfrak{p}) \otimes_{\mathrm{A}} \mathrm{N}\) 满足 \((\mathrm{S}_k)\) （对每个 \(\mathfrak{p} \in \operatorname{Supp}(\mathrm{M})\) ），则 B-模 \(\mathrm{M} \otimes_{\mathrm{A}} \mathrm{N}\) 满足 \((\mathrm{S}_k)\) 。

\begin{enumerate}
\def\labelenumi{\arabic{enumi})}
\setcounter{enumi}{8}
\item
  给出一个局部环 A 和一个有限生成 A-模 M 的例子，使得 \(\mathfrak{m}_{\mathrm{A}}M \neq M\) 且 \(\operatorname{prof}_{A}(M) = +\infty\) （取 A 为一个以无限族未定元的多项式环的局部化）。
\item
  设 A 为 Noether 环，M 为有限生成 A-模，J 为 A 的理想，\((x_{1},\ldots,x_{r})\) 为 J 中元素组成的一个 M-正则序列，且 \(r < \operatorname{prof}_{\mathrm{A}}(\mathrm{J};\mathrm{M})\) 。设 \(q_{1}\ldots,q_{n}\) 为 A 的不含 J 的理想，除可能的两个外均为素理想，并设 \(J_{0}\) 为 J 的一个子集，它对加法和乘法稳定且生成理想 J。证明存在 \(J_{0}\) 的一个元素 x，它不属于任何 \(q_{i}\) ，且使序列 \((x_{1},\ldots,x_{r},x)\) 是 M-正则的（利用 II, § 1, 第 1 节，命题 2 的推论 2）。
\item
  设 \(\mathrm{A}\) 为环，M 为 \(\mathrm{A}\) -模，且 \(x_{1},\ldots ,x_{r}\) 为 \(\mathrm{A}\) 的元素。
\end{enumerate}

\begin{enumerate}
\def\labelenumi{\alph{enumi})}
\item
  对每个满足 \(1 \leqslant p \leqslant r\) 的整数 p，令 \(M_{p} = M/(x_{2}M + \ldots + x_{p}M)\) 。证明：若序列 \((x_{1}, \ldots, x_{r})\) 是 M-正则的，则序列 \((x_{1}, x_{p+1}, x_{p+2}, \ldots, x_{r})\) 是 \(M_{p}\) -正则的（对 p 用归纳法论证）。
\item
  为使序列 \((x_{1}, x_{3})\) 与 \((x_{2}, x_{3})\) 都是 M-正则的，必要且充分的是序列 \((x_{1}x_{2}, x_{3})\) 是 M-正则的。
\item
  设 \(x_p' \in \mathrm{A}\) （ \(1 \leqslant p \leqslant r\) ）。为使序列 \((x_1, \ldots, x_p, \ldots, x_r)\) 与 \((x_1, \ldots, x_p', \ldots, x_r)\) 都是 M-正则的，必要且充分的是序列 \((x_1, \ldots, x_p x_p', \ldots, x_r)\) 是 M-正则的（利用 a) 与 b)）。
\end{enumerate}

\(d\) 设 \(n_1,\ldots ,n_r\) 为 \(\geqslant 1\) 的整数。为使序列 \((x_{1}^{n_{1}},\dots,x_{r}^{n_{r}})\) 是 M-正则的，必要且充分的是序列 \((x_{1},\ldots ,x_{r})\) 是 M-正则的。

\begin{enumerate}
\def\labelenumi{\arabic{enumi})}
\setcounter{enumi}{11}
\tightlist
\item
  设 A 为环，M 为 A-模。
\end{enumerate}

\begin{enumerate}
\def\labelenumi{\alph{enumi})}
\item
  设 \(P \in A[X]\) 。证明：若 M{[}X{]}中以 P 为比的乘法映射的核不为零，则它含有 M 的一个非零元素（设 Q 为 M{[}X{]}中使 PQ = 0 的一个次数最小的元素；若 \(P = a_{0}X^{p} + \ldots + a_{p}\) ，\(Q = m_{0}X^{q} + \ldots + m_{q}\) ，注意 \(a_{0}Q = 0\) ，然后用归纳法注意对所有 i 有 \(a_{i}Q = 0\) ）。
\item
  设 J 为 A 的一个有限生成理想。证明关系 \(\operatorname{prof}_{\mathrm{A}}(\mathrm{J};\mathrm{M}) > 0\) 等价于存在一个 M{[}X{]} -正则的元素（在 JA{[}X{]}中）。
\item
  设 I 与 J 为 A 的有限生成理想。证明等式 \(\mathrm{prof}_{\mathrm{A}}(\mathrm{IJ};\mathrm{M}) = \min (\mathrm{prof}_{\mathrm{A}}(\mathrm{I};\mathrm{M}),\mathrm{prof}_{\mathrm{A}}(\mathrm{J};\mathrm{M}))\) （对 \(\mathrm{prof}_{\mathrm{A}}(\mathrm{IJ};\mathrm{M})\) 用归纳法推理，并利用 \(b)\) ）。
\end{enumerate}

设 A 为环，J 为 A 的理想，M 为 A-模。对每个整数 \(n \geqslant 0\) ，用 \(\operatorname{prof}_{n}(\mathbf{J}; \mathbf{M})\) 表示诸序列长度的上确界（在 \(\overline{\mathbf{N}}\) 中），这些序列由 \(\mathrm{M}[\mathrm{X}_{1}, \ldots, \mathrm{X}_{n}]\) -正则的、属于 \(\mathrm{JA}[\mathrm{X}_{1}, \ldots, \mathrm{X}_{n}]\) 的元素组成。

\begin{enumerate}
\def\labelenumi{\alph{enumi})}
\item
  证明序列 \(\left(\operatorname{prof}_{n}(\mathrm{J};\mathrm{M})\right)_{n\geqslant0}\) 是递增的；用 \(\operatorname{prof}_{\infty}(\mathrm{J};\mathrm{M})\) 表示它的极限（在 \(\overline{\mathbf{N}}\) 中）。证明我们有 \(\operatorname{prof}_{\infty}(\mathrm{J};\mathrm{M})\leqslant\operatorname{prof}_{\mathrm{A}}(\mathrm{J};\mathrm{M})\) 。
\item
  当理想 J 是有限生成的时，证明等式 \(\operatorname{prof}_{\infty}(\mathrm{J};\mathrm{M})=\operatorname{prof}_{\mathrm{A}}(\mathrm{J};\mathrm{M})\) （按第 4 节定理 2 的证明方式推理，并利用习题 12。）
\item
  在一般情形，证明 \(\operatorname{prof}_{\infty}(J;M)\) 是诸数 \(\operatorname{prof}_{\mathrm{A}}(J';M)\) 的上确界，其中 \(J'\) 取遍含于 J 的有限生成理想组成的集合。
\item
  设 \(J'\) 为 A 的一个理想，使得 \(V(J') = V(J)\) 。证明我们有 \(\text{prof}_{\infty}(J'; M) = \text{prof}_{\infty}(J; M)\) 。
\item
  设 A 为一个高为 1 的赋值环，且不是 Noether 的，并设 m 为其极大理想。证明我们有 \(\operatorname{prof}_{\infty}(\mathfrak{m};\mathrm{A})=1\) 和 \(\operatorname{prof}_{\mathrm{A}}(\mathfrak{m};\mathrm{A})\geqslant2\) （参见习题 2）。
\end{enumerate}

\begin{enumerate}
\def\labelenumi{\arabic{enumi})}
\setcounter{enumi}{13}
\tightlist
\item
  设 A 为一个完备 Noether 局部环。
\end{enumerate}

\begin{enumerate}
\def\labelenumi{\alph{enumi})}
\item
  证明：若 A 的一个理想含于 A 的一个由素理想组成的序列之并，则它含于其中的一个（注意拓扑空间 A 是一个 Baire 空间，参见 TG, IX, 第 55 页，定理 1）。
\item
  设 M 为一个具有可数生成族的 A-模。证明定理 2 的结论仍然成立（注意集合 Ass(M) 是可数的，并由 a) 推出：若 \(\mathrm{prof}_{\mathrm{A}}(\mathrm{J};\mathrm{M}) > 0\) ，则存在 J 的一个元素 \(x\) ，使得乘法映射 \(x_{\mathrm{M}}\) 是单射）。
\end{enumerate}

\begin{enumerate}
\def\labelenumi{\arabic{enumi})}
\setcounter{enumi}{14}
\tightlist
\item
  设 A 为 Noether 局部环，并设 M 为一个有限生成 A-模。证明不等式
\end{enumerate}

\[
\operatorname{prof} (A) \leqslant \operatorname{grade} (M) + \dim (M) \leqslant \dim (A).
\]

（对整数 \(g = \text{grade}(M)\) 用归纳法证明第一个不等式，利用第 7 节命题 13 的推论 2 处理 \(g = 0\) 的情形，然后考虑 Ann(M) 的一个 A-正则元素。至于第二个不等式，注意有 grade(M) ≤ prof(A \(_{p}\) )（对每个 p ∈ Supp(M)）。）

\begin{enumerate}
\def\labelenumi{\arabic{enumi})}
\setcounter{enumi}{15}
\tightlist
\item
  设 A 为 Noether 局部环。对每个有限生成 A-模 M，用 \(t_{\mathrm{A}}(\mathrm{M})\) 或简称 \(t(\mathrm{M})\) 表示 \(\kappa_{A}\) -向量空间 \(\operatorname{Ext}_{\mathrm{A}}^{p}(\kappa_{\mathrm{A}},\mathrm{M})\) 的维数，其中 \(p=\operatorname{prof}(\mathrm{M})\) 。
\end{enumerate}

\begin{enumerate}
\def\labelenumi{\alph{enumi})}
\item
  设 \((x_{1},\ldots,x_{r})\) 为 \(m_{A}\) 中对 M 完全割的元素组成的一个序列，它生成一个理想 J。我们有 \(t(\mathrm{M})=t(\mathrm{M}/\mathrm{JM})\) 。
\item
  设 \(\rho : A \to B\) 为 Noether 局部环之间的一个局部同态；设 M 为一个有限生成 A-模，并设 N 为一个在 A 上平坦的有限生成 B-模。证明公式 \(t_{B}(M \otimes_{A} N) = t_{A}(M)t_{B}(\kappa_{A} \otimes_{A} N)\) （利用 \(a\) 化为 \(\mathrm{prof}_{\mathrm{B}}(M \otimes_{A} N) = \mathrm{prof}_{\mathrm{A}}(M) = \mathrm{prof}_{\mathrm{B}}(\kappa_{\mathrm{A}} \otimes_{\mathrm{A}} N) = 0\) 的情形）。
\end{enumerate}

\begin{enumerate}
\def\labelenumi{\arabic{enumi})}
\setcounter{enumi}{16}
\tightlist
\item
  设 A 为 Noether 环，M 为有限生成 A-模。证明 M 是自反的，当且仅当下列两个条件被满足：
\end{enumerate}

\begin{enumerate}
\def\labelenumi{(\roman{enumi})}
\item
  对每个 \(\mathfrak{p} \in \operatorname{Spec}(A)\) （满足 \(\operatorname{prof}(A_{\mathfrak{p}}) \leqslant 1\) ），\(A_{\mathfrak{p}}\) -模 \(M_{\mathfrak{p}}\) 是自反的；
\item
  对每个 \(\mathfrak{p} \in \operatorname{Spec}(\mathrm{A})\) （满足 \(\operatorname{prof}(A_{\mathfrak{p}}) \geqslant 2\) ），有 \(\operatorname{prof}_{\mathrm{A}_{\mathfrak{p}}} (M_{\mathfrak{p}}) \geqslant 2\) 。
\end{enumerate}

（在条件 (i) 与 (ii) 之下，证明典范同态 \(M \to M^{**}\) 的核、然后其余核，都没有相伴素理想。）

\begin{enumerate}
\def\labelenumi{\arabic{enumi})}
\setcounter{enumi}{17}
\tightlist
\item
  设 A 为 Noether 环，J 为 A 的理想，M 与 N 为有限生成 A-模，且 \(\operatorname{prof}_{\mathrm{A}}(J;N)\geqslant2\) 。
\end{enumerate}

\begin{enumerate}
\def\labelenumi{\alph{enumi})}
\item
  证明不等式 \(\mathrm{prof}_{\mathrm{A}}(\mathrm{J};\mathrm{Hom}_{\mathrm{A}}(\mathrm{M},\mathrm{N}))\geqslant 2\) （考虑 J 中的一个 N-正则序列 \((x,y)\) ）。
\item
  假设 \(\mathrm{prof}_{\mathrm{A}}(\mathrm{J};\mathrm{Hom}_{\mathrm{A}}(\mathrm{M},\mathrm{N}))\geqslant 3\) 。证明我们有 \(\mathrm{prof}_{\mathrm{A}}(\mathrm{J};\mathrm{Ext}_{\mathrm{A}}^{1}(\mathrm{M},\mathrm{N}))\geqslant 1\) （考虑与 J 的一个 N-正则元素相伴随的扩张模的正合列）。
\end{enumerate}

\exercisesection{Macaulay 模与环}

\begin{enumerate}
\def\labelenumi{\arabic{enumi})}
\item
  设 A 为 Noether 局部环，并设 \(0 \to M' \to M \to M'' \to 0\) 为 Macaulay A-模的一个正合列。证明我们有 \(\dim(M') = \dim(M)\) ，且 \(\dim(M'')\) 等于 \(\dim(M)\) 或 \(\dim(M) - 1\) 。
\item
  设 A 为 Macaulay 局部环，并设 \(\mathfrak{p} \in \operatorname{Spec}(A)\) 。证明我们有 \(\dim(\widehat{A}/\mathfrak{q}) = \dim(A/\mathfrak{p})\) （对每个 \(\mathfrak{q} \in \operatorname{Ass}(\widehat{A}/\mathfrak{p}\widehat{A})\) ）。特别地，环 \(\widehat{\mathrm{A}}/\mathfrak{p}\widehat{A}\) 没有嵌入素理想。
\item
  \begin{enumerate}
  \def\labelenumii{\alph{enumii})}
  \tightlist
  \item
    设 R 为一个 Noether 的、整的、整闭的 Q-代数，A 为环，并设 \(\rho: \mathrm{R} \to \mathrm{A}\) 为使 A 成为有限 R-代数的一个单射同态。证明子 R-模 \(\rho(\mathrm{R})\) 是 A 的一个直和项（化为命题 8 的推论 3）。
  \end{enumerate}
\end{enumerate}

\begin{enumerate}
\def\labelenumi{\alph{enumi})}
\setcounter{enumi}{1}
\tightlist
\item
  设 A 为一个 Noether 局部 Q-代数。证明 A 满足下列性质（“单项式条件”）：
\end{enumerate}

\begin{enumerate}
\def\labelenumi{(\Roman{enumi})}
\setcounter{enumi}{899}
\tightlist
\item
  对每个极大割序列 \((x_{1},\ldots,x_{d})\) （由 \(m_{A}\) 中的元素组成）以及每个整数 p，元素 \((x_{1}\ldots x_{d})^{p}\) 不属于理想 \((x_{1}^{p+1},\ldots,x_{d}^{p+1})\) 。（化为 A 是完备的情形，从而它具有一个系数域 K，并对一个同态 \(\rho:K[[X_{1},\ldots,X_{d}]]\to A\) （它满足 \(\rho(X_{i})=x_{i}\) ）应用 a）。）
\end{enumerate}

\begin{enumerate}
\def\labelenumi{\arabic{enumi})}
\setcounter{enumi}{3}
\tightlist
\item
  设 A 为局部环，p 为 A 的一个素理想。用 B 表示分式环 \(A[Z][S^{-1}]\) ，其中 S 是形如 \(Z + a\) 的元素（在 A{[}Z{]}中）组成的集合（\(a \in \mathfrak{m}_{A} - \mathfrak{p}\) ）。用 z 表示 B 中的元素 Z/1。
\end{enumerate}

\begin{enumerate}
\def\labelenumi{\alph{enumi})}
\item
  证明 \(\mathrm{B} / (z)\) 同构于 \(\mathbf{A}_{\mathfrak{p}}\) ，且 \(\mathrm{B} / (z - 1)\) 同构于 \(\mathbf{A}\) 。
\item
  以下假设环 A 是整的且 Noether 的，并假设 p 由 A 的一个非零元素 p 生成。设 C 为环 B{[}T{]}/(z(pT-1))；若 A 是 Macaulay 环，则 C 也是。证明 Spec C 有两个不可约分支 \(X_{1}\) 与 \(X_{2}\) ，它们相交于一个闭点 x，使得 \(\text{codim}(\{x\}, X_{1}) = \text{codim}(\{x\}, X_{2}) = 1\) ，\(\text{dim}(X_{1}) = 2\) 且 \(\text{dim}(X_{2}) \geqslant \text{dim}(A_{\mathfrak{p}})\) 。对每个整数 n，给出一个满足 \(\text{dim}(A_{\mathfrak{p}}) = n\) 的例子。
\end{enumerate}

¶ 5) 设 A 为 Noether 局部环，M 为有限生成 A-模，\(\mathbf{x} = (x_{1},\ldots ,x_{n})\) 为 \(\mathfrak{m}_{\mathrm{A}}\) 中元素组成的一个序列，它生成一个理想 \(\mathfrak{x}\)；假设 A-模 M/xM 有有限长度（这蕴含 \(n\geqslant \dim (\mathrm{M})\) ）。用 \(\mathrm{gr(A)}\) 与 \(\mathrm{gr(M)}\) 表示 A 与 M 关于 x-adic 拓扑的伴随分次环，用 \(\xi_{i}\) （ \(1\leqslant i\leqslant n\) ）表示 \(x_{i}\) 在 \(\mathfrak{x} / \mathfrak{x}^{2}\) 中的像，用 \(\xi\) 表示序列 \((\xi_1,\dots ,\xi_n)\) 。

\begin{enumerate}
\def\labelenumi{\alph{enumi})}
\item
  对 \(p \in \mathbf{Z}\) ，令 \(\mathrm{F}^p\mathbf{K}_i(\mathbf{x},\mathrm{M}) = \mathbf{K}_i(\mathbf{x},\mathfrak{x}^{p - i}\mathrm{M})\) （约定 \(\mathfrak{x}^s = A\) ，对 \(s \leqslant 0\) ）。证明 \(\mathrm{F}^p\mathbf{K}_{\bullet}(\mathbf{x},\mathrm{M})\) 是 \(\mathbf{K}_{\bullet}(\mathbf{x},\mathrm{M})\) 的一个子复形，且复形 \(\bigoplus_{n}\mathrm{F}^p\mathbf{K}_{\bullet}(\mathbf{x},\mathrm{M}) / \mathrm{F}^{p + 1}\mathbf{K}_{\bullet}(\mathbf{x},\mathrm{M})\) 等同于 \(\mathbf{K}_{\bullet}^{\mathrm{gr}(A)}(\boldsymbol {\xi},\mathrm{gr}(M))\) 。
\item
  对每个使 H(C) 具有有限长度的 A-模复形 C，令 \(\chi(C) = \sum_{i}(-1)^{i}\lg(\mathrm{H}_{i}(C))\) 。证明存在一个整数 \(p_0\) ，使得复形 \(\mathrm{F}^{p}\mathbf{K}_{\bullet}(\mathbf{x},\mathrm{M}) / \mathrm{F}^{p + 1}\mathbf{K}_{\bullet}(\mathbf{x},\mathrm{M})\) 的同调为零（对 \(p\geqslant p_0\) ）；由此导出等式
\end{enumerate}

\[
\chi \left(\mathbf {K} _ {\bullet} ^ {\operatorname{gr} (\mathrm{A})} (\boldsymbol {\xi}, \operatorname{gr} (\mathrm{M}))\right) = \chi \left(\mathbf {K} _ {\bullet} (\mathrm{x}, \mathrm{M}) / \mathrm{F} ^ {p} \mathbf {K} _ {\bullet} (\mathrm{x}, \mathrm{M})\right) = \sum (- 1) ^ {i} \binom {n} {i} \lg (\mathrm{M} / x ^ {p - i} \mathrm{M})
\]

对 \(p \geqslant p_0\) 成立。

\begin{enumerate}
\def\labelenumi{\alph{enumi})}
\setcounter{enumi}{2}
\item
  证明上列表达式当 \(n > \dim(\mathrm{M})\) 时为零，且等于 \(e_{\mathfrak{x}}(\mathrm{M})\) （当 \(n = \dim(\mathrm{M})\) 时）。
\item
  设 \(h \geqslant 0\) 且 \(p \geqslant p_{0}\)；证明从 \(\mathrm{F}^{p}\mathbf{K}_{\bullet}(\mathbf{x},\mathrm{M})\) 到 \(\mathrm{F}^{p+h}\mathbf{K}_{\bullet}(\mathbf{x},\mathrm{M})\) 的典范单射是一个同调同构。由此推出 \(\mathrm{H}_{i}(\mathrm{F}^{p}\mathbf{K}_{\bullet}(\mathbf{x},\mathrm{M}))\) 对 x-adic 拓扑是离散的。进而证明这个模的由映射 \(\mathrm{H}_{i}(\mathrm{F}^{p+h}\mathbf{K}_{\bullet}(\mathbf{x},\mathrm{M})) \to \mathrm{H}_{i}(\mathrm{F}^{p}\mathbf{K}_{\bullet}(\mathbf{x},\mathrm{M}))\) 的像所定义的滤链是 x-好的，并由此推出 \(\mathrm{F}^{p}\mathbf{K}_{\bullet}(\mathbf{x},\mathrm{M})\) 的同调为零。最后推出 \(\chi(\mathbf{K}_{\bullet}(\mathbf{x},\mathrm{M}))\) 当 \(n > \dim(\mathrm{M})\) 时为零，且等于 \(e_{\mathfrak{x}}(\mathrm{M})\) （当 \(n = \dim(\mathrm{M})\) 时）。
\end{enumerate}

\begin{enumerate}
\def\labelenumi{\arabic{enumi})}
\setcounter{enumi}{5}
\tightlist
\item
  我们保留上一习题的假设；对每个使 N/xN 具有有限长度的有限生成 A-模 N，用 \(\chi(\mathbf{x},\mathbf{N})\) 表示整数 \(\chi(\mathbf{K}_{\bullet}(\mathbf{x},\mathbf{N}))\) 。
\end{enumerate}

\begin{enumerate}
\def\labelenumi{\alph{enumi})}
\item
  对每个子模 \(M'\) （属于 \(M\) ），有 \(\chi(\mathbf{x}, M) = \chi(\mathbf{x}, M') + \chi(\mathbf{x}, M/M')\) 。
\item
  若 \(x_{1}\mathrm{M} = 0\)，证明 \(\chi (\mathbf{x},\mathrm{M})\) 为零。
\item
  用 \(\mathbf{x}'\) 表示序列 \((x_2,\ldots ,x_n)\)；证明等式
\end{enumerate}

\[
\chi (\mathbf {x}, \mathrm{M}) = \chi (\mathbf {x} ^ {\prime}, \mathrm{M} / x _ {1} \mathrm{M}) - \chi (\mathbf {x} ^ {\prime}, \operatorname{Ker} (x _ {1}) _ {\mathrm{M}})
\]

（利用 A, X, 第 207 页，习题 8）。

\begin{enumerate}
\def\labelenumi{\alph{enumi})}
\setcounter{enumi}{3}
\tightlist
\item
  对 \(1 \leqslant i \leqslant n\) ，用 \(\mathrm{K}_i\) 表示以 \(x_i\) 为比的乘法映射（在 \(\mathrm{M}/(x_1\mathrm{M} + \ldots + x_{i-1}\mathrm{M})\) 中的）的核。证明等式
\end{enumerate}

\[
\lg (\mathrm{M} / \mathfrak {x} \mathrm{M}) - \chi (\mathbf {x}, \mathrm{M}) = \sum_ {i = 1} ^ {n} \chi ((x _ {i + 1}, \dots , x _ {n}), \mathrm{K} _ {i}).
\]

\begin{enumerate}
\def\labelenumi{\alph{enumi})}
\setcounter{enumi}{4}
\tightlist
\item
  设 \(p_1, \ldots, p_n\) 为 \(> 1\) 的整数，\(\mathbf{x}^{\mathbf{p}}\) 为序列 \((x_{1}^{p_1}, \ldots, x_n^{p_n})\) 。证明等式 \(\chi(\mathbf{x}^{\mathbf{p}}, \mathrm{M}) = p_1 \ldots p_n \chi(\mathbf{x}, \mathrm{M})\) （先处理 \(n = 1\) 的情形，然后利用 A, X, 第 157 页，推论 2 化为这一情形）。
\end{enumerate}

¶ 7) 设 A 为 Noether 局部环，M 为有限生成 A-模。我们称序列 \((x_1, \ldots, x_n)\) 对 M 是弱正则的，如果对 \(i = 1, \ldots, n\) ，以 \(x_{i+1}\) 为比的乘法映射（在 \(\mathrm{M}/(x_1\mathrm{M} + \ldots + x_i\mathrm{M})\) 中的）的核被 \(\mathfrak{m}_A\) 零化。

\begin{enumerate}
\def\labelenumi{\alph{enumi})}
\item
  若 \(n \leqslant \dim(M)\)，证明对 M 弱正则的序列 \((x_{1}, \ldots, x_{n})\) 对 M 是割的（对 n 用归纳法推理）。给出一个弱正则但不割的序列的例子（可取 \(\mathrm{A} = M = k[X]/(X^{2})\) （其中 k 是一个域）。
\item
  我们称 M 为 Buchsbaum 模，如果对 M 割的每个序列都对 M 弱正则。Macaulay 模是 Buchsbaum 模；若 M 是 Buchsbaum 模，则 \(A_{\mathfrak{p}}\) -模 \(M_{\mathfrak{p}}\) （对每个素理想 \(p \neq \mathfrak{m}_{\mathrm{A}}\)）是 Macaulay 的。 c) 设 M 为 Buchsbaum 模，\((x_{1}, \ldots, x_{n})\) 为对 M 割的一个序列；证明 \(\mathrm{M}/(x_{1}\mathrm{M} + \ldots + x_{n}\mathrm{M})\) 是一个 Buchsbaum 模。
\end{enumerate}

\(d\) ) 为使 M 是 Buchsbaum 模，必要且充分的是：对 M 的每个极大割序列 \((x_{1},\ldots ,x_{n})\) ，在 \(\mathrm{M}_n = \mathrm{M} / (x_1\mathrm{M} + \dots +x_{n - 1}\mathrm{M})\) 中有等式 \(\mathrm{Ker}(x_n)_{\mathrm{M}_n} = \mathrm{Ker}(x_n^2)_{\mathrm{M}_n}\) （先处理 \(n = 1\) 的情形，注意此时 \(\mathrm{Ker}(x_1)_{\mathrm{M}}\) 具有有限长度，并对这个长度用归纳法推理。在一般情形，把每个割序列 \((x_{1},\ldots ,x_{i})\) 补全为极大割序列 \((x_{1},\ldots ,x_{n})\) ，并考虑序列 \((x_{1},\ldots ,x_{i - 1},x_{i + 1}^{k},\ldots ,x_{n}^{k},x_{i})\) （对 \(k\geqslant 1\) ）。）

¶ 8) 设 A 为 Noether 局部环，M 为有限生成 A-模。我们打算证明下列两个条件的等价性：

\begin{enumerate}
\def\labelenumi{(\roman{enumi})}
\item
  M 是 Buchsbaum 模；
\item
  存在一个正整数 \(i(M)\) ，使得有 \(\lg(M/qM)-e_{\mathfrak{q}}(M)=i(M)\) （对每个理想 \(\mathfrak{q}\subset\mathfrak{m}_{A}\) ，它使 \(M/qM\) 具有有限长度）。
\end{enumerate}

\begin{enumerate}
\def\labelenumi{\alph{enumi})}
\item
  在条件 (i) 之下，设 x 为对 M 割的一个序列，并用 x 表示它所生成的理想。利用习题 6, d) 的记号，证明等式 \(\lg(\mathrm{M}/x\mathrm{M}) - e_{\mathfrak{x}}(\mathrm{M}) = \lg(\mathrm{K}_{n})\) 。设 \(\mathbf{x}'\) 为对 M 割的另一个序列，并设 \(\mathrm{K}_{n}^{\prime}\) 为相应的模；证明等式 \(\lg(\mathrm{K}_{n}) = \lg(\mathrm{K}_{n}^{\prime})\) （对 \(n\) 用归纳法推理，考虑一个元素 \(t\) （属于 \(\mathfrak{m}_{\mathrm{A}}\) ），使得序列 \((x_{1},\ldots,x_{n-1},t)\) 与 \((x_{1}',\ldots,x_{n-1}',t)\) 都对 M 割）。
\item
  我们假设条件 (ii) 成立。设 \(\mathbf{x} = (x_{1}, \ldots, x_{n})\) 为对 M 的一个极大割序列，并设 \(M_{n} = \mathrm{M}/(x_{1}\mathrm{M} + \ldots + x_{n-1}\mathrm{M})\) ；对序列 x 与 \((x_{1}, \ldots, x_{n-1}, x_{n}^{2})\) 应用习题 6, d) 的公式（并注意到习题 6, e)），证明有 \(\operatorname{Ker}(x_{n})_{\mathrm{M}_{n}} = \operatorname{Ker}(x_{n}^{2})_{\mathrm{M}_{n}}\) 。最后由习题 7, d) 得出结论）。
\item
  为使 M 是 Buchsbaum 模，必要且充分的是 \(\widehat{\mathbf{M}}\) 也是 Buchsbaum 模；此时我们有 \(i(\widehat{\mathbf{M}}) = i(\mathbf{M})\) 。
\item
  设 M 为 Buchsbaum 模，\((x_{1},\ldots,x_{p})\) 为 \(\mathfrak{m}_{A}\) 中元素组成的对 M 割的一个序列，且 \(p<\dim(M)\) ；证明我们有 \(i(\mathrm{M}/(x_{1}\mathrm{M}+\ldots+x_{p}\mathrm{M}))=i(\mathrm{M})\) （利用 a) 中给出的 \(i(\mathrm{M})\) 的表达式）。
\end{enumerate}

\begin{enumerate}
\def\labelenumi{\arabic{enumi})}
\setcounter{enumi}{8}
\tightlist
\item
  设 A 为 Noether 局部环。
\end{enumerate}

\begin{enumerate}
\def\labelenumi{\alph{enumi})}
\item
  设 M 为一个有限生成 Macaulay A-模，N 为 M 的一个包含 \(\mathfrak{m}_{A}M\) 的子模。证明 N 是一个 Buchsbaum 模，且有 \(i(N) = (\dim(M) - 1)[M/N : \kappa_{A}]\) 。
\item
  若 A 是 Macaulay 环，则理想 \(\mathfrak{m}_{A}\) 是一个 Buchsbaum 模，且当 \(\dim(A) \geqslant 2\) 时它不是 Macaulay 的。
\item
  设 \(k\) 为一个域；我们取 A 为 \(k[[X,Y]]\) 的由 \(X^4, X^3Y, XY^3, Y^4\) 生成的子环。有 \(\dim(A) = 2\) ，\(\text{prof}(A) = 1\) ，且 A 是一个 Buchsbaum 环，其 \(i(A) = 1\) （对 M 取 A 的整闭包并应用 \(a\) ）。
\end{enumerate}

¶ 10) a) 设 \(\rho: A \to B\) 为 Noether 环之间的一个单射同态，使得 B 的子 A-模 \(\mathrm{A}1_{B}\) 是一个直因子。证明若 B 是 Macaulay 环，则 A 也是。

（化为 A 是局部的情形，并通过对 \(\mathrm{prof}_{\mathrm{A}}(\mathrm{B})\) 用归纳法推理，化为 \(\mathrm{Ass}(\mathrm{B})\) 含有一个极大理想的情形。通过考虑 B 中 0 的准素分解，证明此时 B 同构于一个乘积，并对 B 的极大理想的个数用归纳法推理得出结论。

\begin{enumerate}
\def\labelenumi{\alph{enumi})}
\setcounter{enumi}{1}
\tightlist
\item
  设 B 为 Macaulay 环，G 为 B 的一个有限自同构群，其阶在 B 中可逆。证明 G 的不变量环是一个 Macaulay 环。
\end{enumerate}

¶ 11) 设 B 为一个 Noether 整闭环，G 为 B 的一个有限自同构群，A 为 G 的不变量环。

\begin{enumerate}
\def\labelenumi{\alph{enumi})}
\tightlist
\item
  我们假设：对 \(\mathrm{A}\) 的每个高为 1 的素理想 p 以及 B 的每个位于 p 之上的素理想 q，赋值 \(v_{\mathfrak{q}}\) 相对于 \(v_{\mathfrak{p}}\) 是非分歧的（VI, § 8, 第 1 节）。设 \(i: C(A) \to C(B)\) 为除子类群上的典范同态（VII, § 1, 第 10 节）；证明 \(i\) 的核同构于上同调群 \(H^{1}(G,B^{*})\) （A, X, 第 111 页；设 L 为 B 的分式域；考虑 Z \(^{(G)}\) -模的正合列
\end{enumerate}

\[
0 \rightarrow \mathrm{B} ^ {*} \rightarrow \mathrm{L} ^ {*} \rightarrow \mathrm{D} (\mathrm{B}) \rightarrow \mathrm{C} (\mathrm{B}) \rightarrow 0,
\]

并注意 \(\mathrm{H}^{0}(\mathrm{G},\mathrm{D}(\mathrm{B}))\) 等同于 \(\mathrm{D}(\mathrm{A})\) ，且 \(\mathrm{H}^{1}(\mathrm{G},\mathrm{L}^{*})\) 为零）。

\begin{enumerate}
\def\labelenumi{\alph{enumi})}
\setcounter{enumi}{1}
\item
  设 k 为一个特征为 2 的域；取 B 为环 \(k[X_{1},\ldots,X_{4}]\) ，并取 G 为 k-代数 B 的由自同构 \(\sigma\) （它满足 \(\sigma(X_{i})=X_{i+1}\) （对 \(1\leqslant i\leqslant3\) ）且 \(\sigma(X_{4})=X_{1}\) ）所生成的自同构群。由 a) 推出 G 的不变量环 A 是唯一分解的。
\item
  设 m 为 \(\mathrm{A}\) 上理想 \((\mathrm{X}_1, \ldots, \mathrm{X}_4)\) 的迹；令 \(s = \mathrm{X}_1 + \ldots + \mathrm{X}_4\) ，\(t = (\mathrm{X}_1 + \mathrm{X}_3)(\mathrm{X}_2 + \mathrm{X}_4)\) ，\(u = (\mathrm{X}_1 + \mathrm{X}_2)(\mathrm{X}_3 + \mathrm{X}_4)(\mathrm{X}_1 + \mathrm{X}_4)(\mathrm{X}_2 + \mathrm{X}_3)\) ，\(v = \mathrm{X}_1\mathrm{X}_2\mathrm{X}_3\mathrm{X}_4\) 。证明序列 \((s, t, u, v)\) 对 \(\mathrm{A}_{\mathfrak{m}}\) 是割的但不是完全割的。（设 \(x = \mathrm{X}_1^2(\mathrm{X}_2 + \mathrm{X}_3) + \mathrm{X}_2^2(\mathrm{X}_3 + \mathrm{X}_4) + \mathrm{X}_3^2(\mathrm{X}_4 + \mathrm{X}_1) + \mathrm{X}_4^2(\mathrm{X}_1 + \mathrm{X}_2)\) ，\(y = \mathrm{X}_1\mathrm{X}_3 + \mathrm{X}_2\mathrm{X}_4\) ，\(w = \mathrm{X}_1\mathrm{X}_2\mathrm{X}_3 + \mathrm{X}_2\mathrm{X}_3\mathrm{X}_4 + \mathrm{X}_3\mathrm{X}_4\mathrm{X}_1 + \mathrm{X}_4\mathrm{X}_1\mathrm{X}_2\) ；证明我们有 \(sw + y^2 = u\) ，且 \(xy + tw\) 被 \(s\) 整除，从而 \(ux\) 属于理想 \((s, t)\) 。）
\item
  证明 A 不是 Macaulay 环，且子 A-模 \(A.1_{B}\) 不是 B 的直和项。
\end{enumerate}

\exercisesection{深度与同调维数}

\begin{enumerate}
\def\labelenumi{\arabic{enumi})}
\tightlist
\item
  设 \(\rho: A \to B\) 为 Noether 局部环之间的一个局部同态，\(N\) 为一个有限生成 B-模。假设 A-模 \(B\) 与 \(N\) 都是平坦的。
\end{enumerate}

\begin{enumerate}
\def\labelenumi{\alph{enumi})}
\item
  证明等式 \(\mathrm{dp}_{\mathrm{B}}(\mathrm{N}) = \mathrm{dp}_{\kappa_{\mathrm{A}} \otimes_{\mathrm{A}} \mathrm{B}}(\kappa_{\mathrm{A}} \otimes_{\mathrm{A}} \mathrm{N})\) （考虑 N 的一个由有限生成自由 B-模组成的分解）。
\item
  设 M 为一个有限生成 A-模。证明等式
\end{enumerate}

\[
\mathrm{dp} _ {\mathrm{B}} (\mathrm{M} \otimes_ {\mathrm{A}} \mathrm{N}) = \mathrm{dp} _ {\mathrm{A}} (\mathrm{M}) + \mathrm{dp} _ {\mathrm{B}} (\mathrm{N}).
\]

（设 K（相应地，L）为 A-模 M（相应地，B-模 N）的由有限生成自由模组成的分解；证明 \(K \otimes_{A} L\) 是 \(M \otimes_{A} N\) 的一个由自由 B-模组成的分解。若 dp \(_{A}\) (M) = m 且 dp \(_{B}\) (N) = n，计算 \(\mathrm{H}_{m+n}(\kappa_{\mathrm{B}} \otimes_{\mathrm{B}} (\mathrm{K} \otimes_{\mathrm{A}} \mathrm{L}))\) ，并用 \(\mathrm{H}_{m}(\kappa_{\mathrm{A}} \otimes_{\mathrm{A}} \mathrm{K})\) 与 \(\mathrm{H}_{n}(\kappa_{\mathrm{B}} \otimes_{\mathrm{B}} \mathrm{L})\) 来表达。

\begin{enumerate}
\def\labelenumi{\arabic{enumi})}
\setcounter{enumi}{1}
\tightlist
\item
  设 A 为 Noether 环，M 为一个 A-模，它具有一个由有限生成自由模组成的有限自由分解 (L, p)。
\end{enumerate}

\begin{enumerate}
\def\labelenumi{\alph{enumi})}
\tightlist
\item
  证明对 A 的每个素理想 p，有
\end{enumerate}

\[
\sum_ {i \geqslant 0} (- 1) ^ {i} \operatorname{rg} _ {\mathrm{A}} (\mathrm{L} _ {i}) = \sum_ {i \geqslant 0} (- 1) ^ {i} \left[ \operatorname{Tor} _ {i} ^ {\mathrm{A}} (\mathrm{M}, \kappa (\mathfrak {p})): \kappa (\mathfrak {p}) \right] = \sum_ {i \geqslant 0} (- 1) ^ {i} \left[ \operatorname{Ext} _ {\mathrm{A}} ^ {i} (\mathrm{M}, \kappa (\mathfrak {p})): \kappa (\mathfrak {p}) \right];
\]

用 \(\chi(\mathrm{M})\) 表示这个整数。

\begin{enumerate}
\def\labelenumi{\alph{enumi})}
\setcounter{enumi}{1}
\item
  证明有 \(\chi(M) \geqslant 0\) （考虑一个与 A 相伴的理想 \(\mathfrak{p}\) ）。
\item
  证明下列条件等价：
\end{enumerate}

\begin{enumerate}
\def\labelenumi{(\roman{enumi})}
\item
  \(\chi (\mathbf{M}) = 0\)
\item
  M 的零化子不化为零；
\item
  M 的零化子含有一个可消去元素。
\end{enumerate}

（若 \(\chi(M) = 0\) ，证明 \(M_{\mathfrak{p}}\) 为零（对每个 \(\mathfrak{p} \in \operatorname{Ass}(A)\) ）；若 \(\chi(M) > 0\) ，类似地证明 Ann(M) 的零化子含有一个可消去元素，这蕴含 Ann(M) = 0。）

\begin{enumerate}
\def\labelenumi{\arabic{enumi})}
\setcounter{enumi}{2}
\tightlist
\item
  设 A 为 Macaulay 局部环。回忆（A, II, 第 186 页，习题 16）：若 A 的一个理想异于 A 且不是严格包含它的两个理想的交，则称它在 A 中是不可约的。证明下列条件等价：
\end{enumerate}

\begin{enumerate}
\def\labelenumi{(\roman{enumi})}
\item
  A 是 Gorenstein 环；
\item
  由一个极大割序列生成的 A 的每个理想都是不可约的；
\item
  在 A 中存在一个生成不可约理想的极大割序列。
\end{enumerate}

（化为 A 的维数为 0 的情形，并考虑 A 的底座 \(\mathrm{Hom}_{\mathrm{A}}(\kappa_{\mathrm{A}},\mathrm{A})\) 。）

\begin{enumerate}
\def\labelenumi{\arabic{enumi})}
\setcounter{enumi}{3}
\tightlist
\item
  设 A 为 Noether 环，M 为一个非零的有限生成 A-模。有 grade(M) ≤ dpA(M)；若等号成立，则称 M 是完满的。
\end{enumerate}

\begin{enumerate}
\def\labelenumi{\alph{enumi})}
\tightlist
\item
  设 d 为整数；证明下列条件等价：
\end{enumerate}

\begin{enumerate}
\def\labelenumi{(\roman{enumi})}
\item
  A-模 M 是投射维数为 d 的完满模；
\item
  有 \(\operatorname{Ext}_{\mathrm{A}}^{p}(\mathrm{M},\mathrm{A}) = 0\) （对 \(p\neq d\) ）
\item
  对 M 的支撑集中的每个素理想（相应地，极大理想）\(\mathfrak{p}\) ，\(\mathrm{A}_{\mathfrak{p}}\) -模 \(\mathrm{M}_{\mathfrak{p}}\) 是投射维数为 \(d\) 的完满模。
\end{enumerate}

（先证明 (i) 与 (ii) 的等价性。）

\begin{enumerate}
\def\labelenumi{\alph{enumi})}
\setcounter{enumi}{1}
\item
  设 J 为 A 的由一个完全割序列生成的理想，并设 m 为一个整数 \(\geqslant 1\) 。证明 A-模 \(A/J^{m}\) 是完满的。
\item
  若 M 是完满的，则它的相伴素理想是使 \(\text{prof}(A_{\mathfrak{p}}) = \text{dp}_{\mathrm{A}}(M)\) 的那些理想 p。
\item
  对每个投射维数为 d 的完满模 N，令 \(\mathrm{N}^{\prime} = \mathrm{Ext}_{\mathrm{A}}^{d}(\mathrm{N}, \mathrm{A})\) 。证明 \(N^{\prime}\) 是完满的，其投射维数为 d，且 \((\mathrm{N}^{\prime})^{\prime}\) 同构于 N（考虑 N 的一个投射分解）。我们有 \(\mathrm{Ass}(\mathrm{N}^{\prime}) = \mathrm{Ass}(\mathrm{N})\) 。
\item
  设 \(x\) 为 A 的根中的一个非零因子，使得乘法映射 \(x_{\mathrm{M}}\) 是单射。为使 A-模 M 是投射维数为 \(d\) 的完满模，必要且充分的是 (A/xA)-模 M/xM 也是完满的。
\item
  投射维数有限的 Macaulay 模是完满的（化为局部情形，并对 dim(M) 用归纳法论证）。
\item
  若 A 是 Macaulay 环，则完满模恰是投射维数有限的那些 Macaulay 模。
\end{enumerate}

\begin{enumerate}
\def\labelenumi{\arabic{enumi})}
\setcounter{enumi}{4}
\tightlist
\item
  \begin{enumerate}
  \def\labelenumii{\alph{enumii})}
  \tightlist
  \item
    设 \(k\) 为一个域，R 为一个有限次的 \(k\) -代数。为使 R 是 Gorenstein 环，必要且充分的是 R-模 \(\mathrm{Hom}_{k}(\mathrm{R}, k)\) 是秩为 1 的自由模（化为 R 是局部的情形，并注意 R-模 \(\mathrm{Hom}_{k}(\mathrm{R}, k)\) 是内射的且含有一个同构于 \(\kappa_{\mathrm{R}}\) 的子模；另一个证明见 § 8, 第 6 节，命题 5 的推论）。
  \end{enumerate}
\end{enumerate}

\begin{enumerate}
\def\labelenumi{\alph{enumi})}
\setcounter{enumi}{1}
\item
  假设上述条件成立，且代数 R 是 N 型分次的，其中 \(R_{0} = k\) 。设 s 为使 \(R_{s} \neq 0\) 的最大整数。证明 \(R_{s}\) 在 k 上的维数为 1，且等于 R 的底座；对每个非零的 k-线性形式 \(\ell\) （在 \(R_{s}\) 上），k-双线性形式 \((x, y) \mapsto \ell(xy)\) （在 \(R_{i} \times R_{s- i}\) 上）对每个 i 诱导从 \(R_{i}\) 到 \(\operatorname{Hom}_{k}(R_{s-i}, k)\) 的一个同构。
\item
  设 A 为 Noether 环，B 为一个作为有限型平坦模的 A-代数。证明下列条件等价：
\end{enumerate}

\begin{enumerate}
\def\labelenumi{(\roman{enumi})}
\item
  B-模 \(\operatorname{Hom}_{\mathrm{A}}(\mathbf{B},\mathbf{A})\) 是秩为 1 的投射模；
\item
  对 A 的每个极大理想（相应地，素理想）\(\mathfrak{p}\) ，环 \(\kappa(\mathfrak{p})\otimes_{\mathrm{A}}\mathrm{B}\) 是 Gorenstein 环。
\end{enumerate}

当 B 是 Gorenstein 环时，这些条件特别地被满足。

¶ 6) 设 A 为 Noether 局部环，M 为一个具有有限投射维数的有限生成 A-模，N 为一个 A-模，使得 \(\operatorname{Tor}_{i}^{\mathrm{A}}(\mathrm{M},\mathrm{N}) = 0\) （对 i \textgreater{} 0）。

\begin{enumerate}
\def\labelenumi{\alph{enumi})}
\tightlist
\item
  对每个整数 n 定义一个典范同构
\end{enumerate}

\[
\operatorname{Ext} _ {\mathrm{A}} ^ {n} (\kappa_ {\mathrm{A}}, \mathrm{M} \otimes_ {\mathrm{A}} \mathrm{N}) \longrightarrow \bigoplus_ {p \geqslant n} \operatorname{Tor} _ {p - n} ^ {\mathrm{A}} (\kappa_ {\mathrm{A}}, \mathrm{M}) \otimes_ {k} \operatorname{Ext} _ {\mathrm{A}} ^ {p} (\kappa_ {\mathrm{A}}, \mathrm{N}),
\]

并证明有 \(\mathrm{di}_{\mathrm{A}}(\mathrm{M}\otimes_{\mathrm{A}}\mathrm{N})\leqslant \mathrm{di}_{\mathrm{A}}(\mathrm{N})\)。

（设 (P,p) 为 M 的一个有限自由分解，(I,e) 为 N 的一个内射分解，并设 C 为复形 \(P \otimes_A I\) 。证明复形

\[
0 \to \mathrm{C} ^ {0} / \mathrm{B} ^ {0} (\mathrm{C}) \to \mathrm{C} ^ {1} \to \mathrm{C} ^ {2} \to \dots
\]

定义 \(M \otimes_A N\) 的一个内射分解。）

\begin{enumerate}
\def\labelenumi{\alph{enumi})}
\setcounter{enumi}{1}
\item
  证明不等式 \(\mathrm{prof}_{\mathrm{A}}(\mathrm{N}) = \mathrm{dp}_{\mathrm{A}}(\mathrm{M}) + \mathrm{prof}(\mathrm{M} \otimes_{\mathrm{A}} \mathrm{N})\) 。由此导出 Auslander-Buchsbaum 定理的另一个证明。
\item
  进而假设 N 具有有限投射维数。证明等式 \(\mathrm{dp}_{\mathrm{A}}(\mathrm{M}\otimes_{\mathrm{A}}\mathrm{N})=\mathrm{dp}_{\mathrm{A}}(\mathrm{M})+\mathrm{dp}_{\mathrm{A}}(\mathrm{N})\) 。
\item
  假设 A 是 Gorenstein 环，并用 d 表示它的维数。设 n 为整数；证明 \(\kappa_{A}\) -向量空间 \(\operatorname{Ext}_{\mathrm{A}}^{n}(\mathbf{k}_{\mathrm{A}},\mathrm{M})\) 、\(\operatorname{Tor}_{d-n}^{\mathrm{A}}(\mathbf{k}_{\mathrm{A}},\mathrm{M})\) 与 \(\operatorname{Ext}_{\mathrm{A}}^{d-n}(\mathrm{M},\mathbf{k}_{\mathrm{A}})\) 具有相同的维数。
\end{enumerate}

\begin{enumerate}
\def\labelenumi{\arabic{enumi})}
\setcounter{enumi}{6}
\tightlist
\item
  设 A 为 Noether 局部环，M 为一个有限生成 A-模，T 为一个具有有限内射维数的有限生成 A-模。证明等式
\end{enumerate}

\[
\operatorname{prof} (\mathrm{A}) = \operatorname{prof} _ {\mathrm{A}} (\mathrm{M}) + \sup \left\{i \mid \operatorname{Ext} _ {\mathrm{A}} ^ {i} (\mathrm{M}, \mathrm{T}) \neq 0 \right\}
\]

（对 \(\mathrm{prof}_{\mathrm{A}}(\mathrm{M})\) 用归纳法推理，并利用命题 9）。

\begin{enumerate}
\def\labelenumi{\arabic{enumi})}
\setcounter{enumi}{7}
\tightlist
\item
  设 A 为环，并设 \(u: \mathrm{A}^p \to \mathrm{A}^q\) 为一个 A-线性映射。我们称 \(u\) 的秩——并用 \(\text{rg}(u)\) 表示——为最大的整数 \(r\) ，它使 \(\wedge^r u \neq 0\) 。用 \(\mathfrak{d}(u)\) 表示由 \(\text{rg}(u)\) 阶子式（属于 \(u\) 的矩阵）所生成的 A 的理想。
\end{enumerate}

\begin{enumerate}
\def\labelenumi{\alph{enumi})}
\item
  证明为使 \(\mathfrak{d}(u)\) 等于 A，必要且充分的是 \(u\) 的余核是秩为 \(q - \mathrm{rg}(u)\) 的投射模（化为 \(\mathrm{A}\) 是局部的情形，并在一个适当的基中写出 \(u\) 的矩阵）。
\item
  设 \(v: \mathbf{A}^q \to \mathrm{A}^s\) 为一个 A-线性映射；我们假设有 \(\mathfrak{d}(u) = \mathfrak{d}(v) = \mathbf{A}\) 为使 \(\operatorname{Ker} v = \operatorname{Im} u\) ，必要且充分的是有 \(\operatorname{rg}(u) + \operatorname{rg}(v) = q\) （同样的方法）。
\item
  若 \(\operatorname{rg}(u) = q\) ，证明 \(\mathfrak{d}(u)\) 零化 \(u\) 的余核（通过化为 \(p = q\) 的情形）。
\end{enumerate}

¶ 9) 设 A 为环，并设 (L, d) 为一个由有限生成自由 A-模组成的有界复形，它在次数 \textless{} 0 处为零。我们打算证明下列条件等价：

\begin{enumerate}
\def\labelenumi{(\roman{enumi})}
\item
  有 \(\mathrm{H}_i(\mathrm{L}) = 0\) （对 \(i > 0\) ）；
\item
  有 \(\operatorname{prof}_{\mathrm{A}}(\mathfrak{d}(d_i); \mathrm{A}) \geqslant i\) 且 \(\operatorname{rg}(d_i) + \operatorname{rg}(d_{i+1}) = \operatorname{rg}_{\mathrm{A}}(\mathrm{L}_i)\) （对每个 \(i \geqslant 1\) ）。
\end{enumerate}

\begin{enumerate}
\def\labelenumi{\alph{enumi})}
\item
  证明只需对环 A 是 Noether 的情形证明这一命题（考虑由矩阵 \(d_{i}\) （对 \(i \geqslant 1\) ）的系数所生成的 A 的子环）。以下假设这个条件成立。
\item
  若 A 是一个维数为 0 的局部环，证明 (ii) 蕴含 (i)（利用习题 8, b)）。
\item
  假设环 A 是局部的，条件 (ii) 成立，且 \(\operatorname{Supp}(\mathrm{H}_i(\mathrm{L})) = \{\mathfrak{m}_{\mathrm{A}}\}\) （对 \(i \geqslant 1\) ）。证明 \(\mathrm{H}_i(\mathrm{L}) = 0\) （对 \(i \geqslant 0\) ）（设 \(p\) 为使 \(\mathfrak{d}(d_p) \neq \mathrm{A}\) 的最大整数；由习题 8, b) 推出 \(\mathrm{H}_i(\mathrm{L}) = 0\) （对 \(i > p\) ），然后由习题 8, a) 以及命题 3 的推论 2 （ \(n^{\circ} 2\) ）推出复形 \(0 \to L_p / B_p(L) \to L_{p-1} \to \ldots \to L_0\) 在次数 \(>0\) 处是零调的）。
\item
  证明蕴含关系 (ii)⇒(i)（化为 A 是局部的情形，并对 dim(A) 用归纳法论证）。
\item
  我们假设复形 L 在次数 \textgreater{} 0 处是零调的。设 p 为与 A 相伴的一个素理想。证明对每个 \(i \geqslant 1\) ，(Coker \(d_{i}\) )p 是一个自由 Ap-模，其秩为 \(\sum_{p \geqslant i-1} (-1)^{p} \operatorname{rg}_{\mathrm{A}}(\mathrm{L}_{p})\) ；此外，我们有 \(\operatorname{rg}((d_{i})_{\mathfrak{p}}) = \sum_{p \geqslant i} (-1)^{p-i} \operatorname{rg}_{\mathrm{A}}(\mathrm{L}_{p})\) 以及 \(\mathfrak{d}((d_{i})_{\mathfrak{p}}) = A_{\mathfrak{p}}\) （习题 8）。
\end{enumerate}

\(f)\) 于是我们有 \(\operatorname{rg}(d_i) = \sum_{p \geqslant i} (-1)^{p - i} \operatorname{rg}_A(L_p)\) ，从而有 \(\operatorname{rg}(d_i) + \operatorname{rg}(d_{i+1}) = \operatorname{rg}_\mathrm{A}(L_i)\) （对每个 \(i \geqslant 1\) ）。

\(g)\) 设 \(i\) 为一个整数 \(\geqslant 1\) ，并设 \(\mathfrak{p}\) 为 \(\mathrm{A}\) 的一个素理想。证明条件 \(\mathrm{dp}_{\mathrm{A}_{\mathfrak{p}}}(H_0(L)_{\mathfrak{p}}) \geqslant i\) 等价于 \(\mathfrak{d}(d_i)A_{\mathfrak{p}} \neq A_{\mathfrak{p}}\) ，亦即 \(\mathfrak{d}(d_i) \subset \mathfrak{p}\) （利用习题 8, a)）。利用 § 1, 第 5 节的命题 8 以及定理 1，推出有 \(\operatorname{prof}_{\mathrm{A}}(\mathfrak{d}(d_i); A) \geqslant i\) ，这就证明了蕴含关系 (i) \(\Rightarrow\) (ii)。

¶ 10) 我们保留习题 9 的假设；我们假设复形 L 在次数 ≠ 0 处是正合的。

\begin{enumerate}
\def\labelenumi{\alph{enumi})}
\item
  证明包含关系 \(V(\mathfrak{d}(d_{i+1})) \subset V(\mathfrak{d}(d_i))\) （对每个 \(i \geqslant 1\) ）（注意由习题 8，素理想 p 不含 \(\mathfrak{d}(d_i)\) 当且仅当 \(A_{\mathfrak{p}}\)-模 \((\operatorname{Coker} d_i)_{\mathfrak{p}}\) 是自由的）。
\item
  我们现在假设 \(\mathrm{H}_{0}(\mathrm{L})\) 的零化子不约化为零。证明有 \(\mathrm{rg}(d_{1}) = \mathrm{rg}_{\mathrm{A}}(\mathrm{L}_{0})\) （若 f 是 \(\mathfrak{o}(d_{1})\) 的一个可消去元素，则 \(A_{f}\) -模 \(\mathrm{H}_{0}(\mathrm{L})_{f}\) 依习题 8 是具有常秩的投射模，故必为零）；由此推出 \(\mathrm{H}_{0}(\mathrm{L})\) 的支撑集是 \(\mathrm{V}(\mathfrak{o}(d_{1}))\) 。
\item
  我们置 \(p = \text{prof}_A(\mathfrak{d}(d_1); A)\) 。证明有 \(V(\mathfrak{d}(d_i)) = \text{Supp}(\mathbf{H}_0(L))\) （对 \(1 \leqslant i \leqslant p\) ）（如 \(a\) ) 的证明中那样推理，并注意：若一个素理想 \(\mathfrak{p}\) （属于 \(\mathrm{H}_0(L)\) 的支撑集）不含 \(\mathfrak{d}(d_k)\) ，则将有 \(\text{dp}_{\mathrm{A}_{\mathfrak{p}}} \, H_0(L)_{\mathfrak{p}} < k \leqslant \text{grade}_{\mathrm{A}_{\mathfrak{p}}} \, H_0(L)_{\mathfrak{p}}\) 。
\end{enumerate}

\begin{enumerate}
\def\labelenumi{\arabic{enumi})}
\setcounter{enumi}{10}
\tightlist
\item
  设 A 为局部环。
\end{enumerate}

\begin{enumerate}
\def\labelenumi{\alph{enumi})}
\tightlist
\item
  设 n 为一个整数 \(\geqslant 1\) ，\(u : A^{n-1} \to A^n\) 为一个 A-线性映射，\(\mathfrak{d}(u)\) 为由 u 的 n-1 阶子式所生成的理想（习题 8）。若 \(\operatorname{prof}_{\mathrm{A}}(\mathfrak{d}(u); A) \geqslant 2\) ，证明 \(\mathfrak{d}(u)\) 容许分解
\end{enumerate}

\[
0 \rightarrow \mathrm{A} ^ {n - 1} \xrightarrow {u} \mathrm{A} ^ {n} \xrightarrow {v} \mathfrak {d} (u) \rightarrow 0,
\]

其中 \(v\) 由同态 \(\wedge^{n-1}(^t u): A^n \to A\) 诱导（应用习题 9）。

\begin{enumerate}
\def\labelenumi{\alph{enumi})}
\setcounter{enumi}{1}
\tightlist
\item
  反之，设 J 为 A 的一个投射维数为 1 的理想。证明存在 A 的一个可消去元素 a、一个整数 n 以及一个 A-线性映射 \(u: A^{n-1} \to A^n\) ，使得 \(\mathrm{J} = a\mathfrak{d}(u)\) ；有 \(\operatorname{prof}_{\mathrm{A}}(\mathrm{J}; \mathrm{A}) = 2\) （“Hilbert-Burch 定理”：由习题 9 推出 J 的一个极小分解形如 \(0 \to A^{n-1} \xrightarrow{u} \mathrm{A}^n \to J\) ，其中 \(\operatorname{prof}_{\mathrm{A}}(\mathfrak{d}(u); \mathrm{A}) \geqslant 2\) ，然后应用 a) 以及 § 1 的习题 1，并注意 \(\mathfrak{d}(u)\) 含有 A 的一个可消去元素）。
\end{enumerate}

\begin{enumerate}
\def\labelenumi{\arabic{enumi})}
\setcounter{enumi}{11}
\tightlist
\item
  设 A 为 Noether 局部环，并设 M 与 N 为非零的有限生成 \(\mathrm{A}\)-模。假设 M 的投射维数有限。设 \(q\) 为使 \(\operatorname{Tor}_q^A(M, N) \neq 0\) 的最大整数。
\end{enumerate}

\begin{enumerate}
\def\labelenumi{\alph{enumi})}
\item
  假设 \(\mathrm{prof}_{\mathrm{A}}(\mathrm{Tor}_{q}^{\mathrm{A}}(\mathrm{M},\mathrm{N})) = 0\) 。证明等式 \(\mathrm{prof}_{\mathrm{A}}(\mathrm{N}) + q = \mathrm{dp}_{\mathrm{A}}(\mathrm{M})\) （可对 \(\mathrm{prof}_{\mathrm{A}}(\mathrm{N})\) 用归纳法推理，考虑一个元素 \(x\) （属于 \(\mathfrak{m}_{\mathrm{A}}\) ，且在 N 中为非零因子））。
\item
  我们进一步假设 A-模 \(M \otimes_{A} N\) 具有有限长度。由 a) 推出：\(\operatorname{Tor}_{i}^{A}(M, N)\) 对 i \textgreater{} depth(A) - depth(M) - depth(N) 为零，且在等号成立时非零。特别地，\(\operatorname{prof}(M) + \operatorname{prof}(N) \leqslant \operatorname{prof}(A)\) 。
\end{enumerate}

¶ 13) 设 A 为 Noether 局部环。假设下列条件被满足 \(^{1}\) ：

（GM） 对每个 Noether 局部 A-代数 B，存在一个 B-模 M，使得 \(\mathfrak{m}_{\mathrm{B}}\mathrm{M}\neq\mathrm{M}\) 且 \(\operatorname{prof}_{\mathrm{B}}(\mathrm{M})=\dim(\mathrm{B})\) 。

\begin{enumerate}
\def\labelenumi{\alph{enumi})}
\item
  设 P 为一个由平坦 A-模组成的、有限长度为 \(\ell\) 的复形，使得 \(\operatorname{Supp}(\mathrm{H}(\mathrm{P})) = \{\mathfrak{m}_{\mathrm{A}}\}\) 。证明有 \(\ell \geqslant \dim(\mathrm{A})\) （利用 § 1 的习题 6） \(^2\) 。
\item
  设 B 为 Noether 环，\(\rho: A \to B\) 为一个同态，\(^a\rho: \operatorname{Spec}(B) \to \operatorname{Spec}(A)\) 为相应的映射，M 为一个有限生成 A-模。证明 \(^a\rho^{-1}(\operatorname{Supp}(M))\) 在 \(\operatorname{Spec}(B)\) 中的余维数 \(\leqslant \mathrm{dp}_A(M)\) （把 \(a\) ) 应用于复形 \(L \otimes_A B_q\) ，其中 \(q\) 是 \(^a\rho^{-1}(\operatorname{Supp}(M))\) 的一个分支的极小素理想，而 L 是 M 的一个自由分解）。
\item
  设 M、N 为有限生成 A-模，其中 \(M \neq 0\) 。证明不等式
\end{enumerate}

\[
\dim (N) \leqslant \mathrm{dp} _ {A} (M) + \dim (M \otimes_ {A} N)
\]

（相交定理：对整数 \(d = \dim(M \otimes_{\mathrm{A}} N)\) 用归纳法推理。若 d = 0，把 b) 应用于环 B = A/Ann(M)；一般情形考虑 A 的一个对 N 以及对 M \(\otimes_{A} N\) 正则的元素）。

\begin{enumerate}
\def\labelenumi{\alph{enumi})}
\setcounter{enumi}{3}
\tightlist
\item
  为使存在一个具有有限长度且有限投射维数的 A-模，必要且充分的是 A 为 Macaulay 环。
\end{enumerate}

\begin{enumerate}
\def\labelenumi{\arabic{enumi})}
\setcounter{enumi}{13}
\tightlist
\item
  设 \(\mathrm{A}\) 为满足条件 (GM) 的 Noether 局部环（习题 13），M 为一个具有有限投射维数的有限生成 A-模。
\end{enumerate}

\begin{enumerate}
\def\labelenumi{\alph{enumi})}
\tightlist
\item
  证明不等式
\end{enumerate}

\[
\dim (A) \leqslant \mathrm{dp} _ {A} (M) + \dim (M) \quad \dim (A) - \operatorname{prof} (A) \leqslant \dim (M) - \operatorname{prof} (M).
\]

若此外模 M 是完满的（习题 4），则有 \(\dim(A) = \mathrm{dp}_{A}(M) + \dim(M)\) （利用 § 1 的习题 15）。

\begin{enumerate}
\def\labelenumi{\alph{enumi})}
\setcounter{enumi}{1}
\item
  设 \(\mathfrak{p} \in \mathrm{Supp}(M)\) ；若 \(\mathrm{A}_{\mathfrak{p}}\) -模 \(\mathrm{M}_{\mathfrak{p}}\) 是 Macaulay 的，则环 \(\mathrm{A}_{\mathfrak{p}}\) 是 Macaulay 环。特别地，当 \(\mathfrak{p}\) 是 \(\mathrm{Supp}(M)\) 的一个极小素理想时即属此情形。
\item
  设 \(\mathbf{x} = (x_{1},\ldots ,x_{r})\) 为 \(\mathfrak{m}_{\mathrm{A}}\) 中元素的一个正则序列，它生成一个具有有限投射维数的理想 J。证明序列 \(\mathbf{x}\) 对 A 是完全正则的（设 \(\mathfrak{p}\in V(J)\) ，并设 \(\mathfrak{q}\subset \mathfrak{p}\) 为 \(V(J)\) 的一个极小素理想；注意有 \(\mathrm{prof}(A_{\mathfrak{p}})\geqslant \mathrm{dp}_{A_{\mathfrak{p}}}(A_{\mathfrak{p}} / J_{\mathfrak{p}})\geqslant \mathrm{dp}_{A_{\mathfrak{q}}}(A_{\mathfrak{q}} / J_{\mathfrak{q}})\) ，且有 \(\mathrm{dp}_{A_{\mathfrak{q}}}(A_{\mathfrak{q}} / J_{\mathfrak{q}})\geqslant r\) （依 \(b\) ）），从而最终有 \(\mathrm{prof}_{\mathrm{A}}(\mathrm{J};\mathrm{A})\geqslant r\) ；借助 § 1, 第 3 节的定理 1 的推论 1，断定 \(\mathbf{x}\) 是完全正则的）。
\item
  证明每个 M-正则序列都是 A-正则的。（证明 A 的每个零化子非零的元素 a 都不是 M-正则的；为此，通过局部化化为 Supp(M) ∩ Supp(a) = \(\{\mathfrak{m}_{A}\}\) 的情形；把习题 13 a) 应用于复形 \(L \otimes_A A/p\)（其中 L 是 M 的一个自由分解，而 p 是 a 的一个相伴素理想），证明不等式 \(\operatorname{prof}(A) \leqslant \dim(A/p) \leqslant \operatorname{dp}_{A}(M)\) ，并由此推出 prof(M) = 0。）
\item
  为使 A 是整环，必要且充分的是它含有一个投射维数 \(< +\infty\) 的素理想（若 \(\mathrm{dp}_{\mathrm{A}}(\mathfrak{p}) < +\infty\) ，证明 A 可等同于正则环 \(\mathrm{A}_{\mathfrak{p}}\) 的一个子环）。
\end{enumerate}

\begin{enumerate}
\def\labelenumi{\arabic{enumi})}
\setcounter{enumi}{14}
\tightlist
\item
  设 A 为 Noether 局部环，M 为一个具有有限内射维数的有限生成 A-模。
\end{enumerate}

\begin{enumerate}
\def\labelenumi{\alph{enumi})}
\item
  设 \(\mathfrak{p} \in \operatorname{Supp}(M)\) 。证明等式 \(\operatorname{prof}(A) = \operatorname{prof}(A_{\mathfrak{p}}) + \dim(A/p)\) 。（设 \(p = \operatorname{prof}(A_{\mathfrak{p}})\) 且 \(q = \dim(A/p)\) ；注意有 \(p = \operatorname{di}_{A_{\mathfrak{p}}} (M_{\mathfrak{p}})\) ，从而 \(\operatorname{Ext}_{A_{\mathfrak{p}}}^{p}(\kappa(\mathfrak{p}), M_{\mathfrak{p}}) \neq 0\) ，故依 § 1, 第 7 节的引理 3 有 \(\operatorname{Ext}_{A}^{p+q}(\kappa_A, M) \neq 0\) ，这蕴含 \(\operatorname{prof}(A) = \operatorname{di}_A(M) = p + q\) 。）若 \(\operatorname{Supp}(M) = \operatorname{Spec}(A)\) ，则 A 是 Macaulay 环。
\item
  证明等式 grade(M) + dim(M) = prof(A)（利用 a) 以及 § 1 的习题 15）。
\end{enumerate}

¶ 16) 设 A 为 Macaulay 环，M 为一个具有有限投射维数的 A-模。证明为使 M 是平坦的，必要且充分的是：对 A 的每个由完全割序列生成的理想 J，典范同态 \(J \otimes_{A} M \to JM\) 是双射（用对序列长度的归纳法证明这个条件蕴含 \(\operatorname{Tor}_{i}^{A}(A/J, M) = 0\) （对一切 i \textgreater{} 0），然后通过对 n 的递降归纳法证明 \(\operatorname{Tor}_{n}^{A}(N, M) = 0\) （对每个有限生成 A-模 N 及每个整数 n \textgreater{} 0））。例子：Dedekind 环上的无挠模是平坦的（参看 VII, § 2, 习题 14）。
